\part{基础数据结构}

\chapter{并查集}

\section{普通并查集}
\entrymeta{$O(\alpha(n))$}
带按秩合并（size）与路径压缩的并查集。下标从 1 开始，\texttt{init(n)} 后每个点自成一个集合。


\begin{codeline}
const int DSU_N = 100005;

class DSUNode {
public:
    int fa;
    int r;
};

class DSU {
public:
    DSUNode node[DSU_N];

    /*
     * 初始化编号1..maxn的独立集合
     */
    void init(int maxn) {
        for(int i = 0; i <= maxn; i++) {
            node[i] = {i, 1};
        }
    }

    /*
     * 查询u所在集合的代表元并压缩路径
     */
    int find(int u) {
        return node[u].fa == u ? u : node[u].fa = find(node[u].fa);
    }

    /*
     * 合并u、v所在集合，返回合并后的代表元
     */
    int join(int u, int v) {
        u = find(u);
        v = find(v);
        if(u != v) {
            if(node[v].r == node[u].r) {
                node[u].fa = v;
                node[v].r++;
            } else {
                if(node[v].r < node[u].r) {
                    swap(u, v);
                }
                node[u].fa = v;
            }
        }
        return v;
\end{codeline}
\vspace{-0.9\baselineskip}
\begin{codeline}[firstnumber=47]
    }
};
\end{codeline}
\begin{warnbox}
下标从 1 开始；\texttt{init(n)} 后每个点自成一个集合。
\end{warnbox}

\section{可持久化并查集}
\entrymeta{$O(\log^2 n)$}
\subsection{P3402 【模板】可持久化并查集}\label{p3402-ux6a21ux677fux53efux6301ux4e45ux5316ux5e76ux67e5ux96c6}


\begin{codeline}
class Data {
public:
    int fa = 0, dep = 0;
};

class PersistentDSU {
public:
    PersistentArray<Data> arr;
    vector<int> version;
    /*
     * 查询版本ver中x所在集合的根信息
     */
    Data find(int ver, int x) const {
        Data d = arr.query(version[ver], x);
        while(d.fa != x) {
            x = d.fa;
            d = arr.query(version[ver], x);
        }
        return d;
    }
    /*
     * 建立n个独立集合，初始版本为0
     */
    void init(int n, size_t capacity = 0) {
        vector<Data> a(n + 1);
        for(int i = 1; i <= n; i++) {
            a[i].fa = i;
        }
        arr.build(a, capacity);
        version.assign(1, 0);
    }
    /*
     * 复制版本ver，返回新版本号
     */
    int copy(int ver) {
        version.push_back(version[ver]);
        return int(version.size()) - 1;
    }
    /*
     * 基于版本ver合并x、y所在集合，返回新版本号
     */
    int merge(int ver, int x, int y) {
        Data a = find(ver, x), b = find(ver, y);
        int root = version[ver];
        if(a.fa != b.fa) {
            if(a.dep > b.dep) {
\end{codeline}
\vspace{-0.9\baselineskip}
\begin{codeline}[firstnumber=47]
                swap(a, b);
            }
            root = arr.modify(root, a.fa, {b.fa, a.dep});
            if(a.dep == b.dep) {
                root = arr.modify(root, b.fa, {b.fa, b.dep + 1});
            }
        }
        version.push_back(root);
        return int(version.size()) - 1;
    }
    /*
     * 判断版本ver中x、y是否属于同一集合
     */
    bool same(int ver, int x, int y) const {
        return find(ver, x).fa == find(ver, y).fa;
    }
};
\end{codeline}
\begin{warnbox}
按秩合并、不路径压缩；回滚时把根指向旧版本根。
\end{warnbox}

\section{区间并查集}
\entrymeta{$O((n+m)\log n\,\alpha(n))$}
\begin{codeline}
const int64 N = 100005;
const int64 MODER = 1000000007;
const int64 LOG_N = 20;
int64 n, m;
int64 lg[N];

class Node {
public:
    int64 fa;
    int64 w;
};

Node f[N][LOG_N];

/*
 * 返回以idx为起点、长度为2^log_gap的区间代表元
 */
int find_root(int64 idx, int64 log_gap) {
    int root = idx;
    while(f[root][log_gap].fa != -1) {
        root = f[root][log_gap].fa;
    }
    while(idx != root) {
        int fa = f[idx][log_gap].fa;
        f[idx][log_gap].fa = root;
        idx = fa;
    }
    return root;
}

/*
 * 合并u、v为起点的两个等长区间
 */
void cat(int64 u, int64 v, int64 log_gap) {
    if(u == -1 || v == -1 || u > n || v > n) {
        return;
    }

    u = find_root(u, log_gap);
    v = find_root(v, log_gap);
    if(u != v) {
        if(f[u][log_gap].w > f[v][log_gap].w) {
            swap(u, v);
        }
        f[u][log_gap].fa = v;
        f[v][log_gap].w += f[u][log_gap].w;
\end{codeline}
\vspace{-0.9\baselineskip}
\begin{codeline}[firstnumber=47]
    }
}

void push_down(int idx, int log_gap) {
    if(log_gap > 0 && f[idx][log_gap].fa != -1) {
        cat(idx, f[idx][log_gap].fa, log_gap - 1);
        cat(idx + (1 << (log_gap - 1)), f[idx][log_gap].fa + (1 << (log_gap - 1)), log_gap - 1);
        f[idx][log_gap].fa = -1;
    }
}

int main() {
    ios::sync_with_stdio(false);
    cin.tie(nullptr);

    cin >> n >> m;

    for(int i = 2; i <= n; i++) {
        lg[i] = lg[i >> 1] + 1;
    }

    for(int i = 0; i <= n; i++) {
        for(int j = 0; j <= lg[n]; j++) {
            f[i][j].fa = -1;
            f[i][j].w = 1;
        }
    }

    int64 l, r, L, R, lag, gap;
    while(m--) {
        cin >> l >> r >> L >> R;

        lag = r - l + 1;
        gap = 1 << (lg[n] + 1);
        while(gap) {
            if(gap & lag) {
                cat(l + ((lag - lag % gap) ^ gap), L + ((lag - lag % gap) ^ gap), lg[gap]);
            }
            gap >>= 1;
        }
    }
    for(int j = lg[n]; j >= 0; j--) {
        for(int i = 1; i + (1ll << j) - 1 <= n; i++) {
            if(f[i][j].fa != -1) {
                push_down(i, j);
            }
\end{codeline}
\vspace{-0.9\baselineskip}
\begin{codeline}[firstnumber=93]
        }
    }
    int cnt = 0;

    for(int i = 1; i <= n; i++) {
        if(f[i][0].fa == -1) {
            cnt++;
        }
    }

    int64 res = 9;
    cnt--;
    for(int i = 0; i < cnt; i++) {
        res = res * 10 % MODER;
    }
    cout << (res) << '\n';
}
\end{codeline}
\begin{warnbox}
倍增并查集：同一层内合并，向下传播。
\end{warnbox}


\chapter{树状数组}
\entrymeta{$O(\log n)$}
\begin{codeline}
template <class Data = long long>
class BIT {
public:
    vector<Data> arr;
    int size = 0;
    /*
     * 初始化长度为n的树状数组
     */
    void init(int n) {
        size = n;
        arr.assign(n + 1, Data());
    }
    /*
     * 在idx处加x，下标从1开始
     */
    void add(int idx, Data x) {
        for(int i = idx; i <= size; i += lowBit(i)) {
            arr[i] += x;
        }
    }
    /*
     * 查询前缀[1,r]的和，r=0时返回0
     */
    Data prefix(int r) const {
        Data res{};
        for(int i = r; i; i -= lowBit(i)) {
            res += arr[i];
        }
        return res;
    }
};
\end{codeline}
\begin{warnbox}
下标从 1 开始；\texttt{prefix(0)} 返回幺元。
\end{warnbox}


\chapter{线段树}

\section{一般线段树}
\entrymeta{$O(\log n)$}
\begin{codeline}

class Data {
public:
    int64 sum = 0;
    Data operator+(const Data& o) const {
        return {sum + o.sum};
    }
};

class Tag {
public:
    int64 add = 0, mul = 1;
    // a + b：先执行 a，再执行 b。
    Tag operator+(const Tag& o) const {
        return {add * o.mul + o.add,
                mul * o.mul};
    }
    Data apply(const Data& d, int l, int r) const {
        return {d.sum * mul + add * (r - l + 1)};
    }
};
\end{codeline}


\subsection{通用实现}\label{ux901aux7528ux5b9eux73b0}

下面的类不依赖上述示例；使用时提供自己的 \texttt{Data\ /\ Tag}，实例化 \texttt{SegmentTree\textless{}Data,\ Tag\textgreater{}}。


\begin{codeline}
template <class Value, class Tag>
class SegmentTree {
public:
    class Node {
    public:
        Value data{};
        Tag tag{};
    };
    int n = 0;
    vector<Node> node;
    void push_up(int u) {
        node[u].data = node[u * 2].data + node[u * 2 + 1].data;
    }
    void apply(int u, int l, int r, const Tag& t) {
        node[u].data = t.apply(node[u].data, l, r);
        node[u].tag = node[u].tag + t;
    }
    void push_down(int u, int l, int r) {
        int mid = l + (r - l) / 2;
        apply(u * 2, l, mid, node[u].tag);
        apply(u * 2 + 1, mid + 1, r, node[u].tag);
        node[u].tag = Tag();
    }
    void build(int u, int l, int r) {
        if(l == r) {
            node[u].data = Value{};
            return;
        }
        int mid = l + (r - l) / 2;
        build(u * 2, l, mid);
        build(u * 2 + 1, mid + 1, r);
        push_up(u);
    }
    void mod(int u, int l, int r, int ql, int qr, const Tag& t) {
        if(ql <= l && r <= qr) {
            apply(u, l, r, t);
            return;
        }
        int mid = l + (r - l) / 2;
        push_down(u, l, r);
        if(ql <= mid) {
            mod(u * 2, l, mid, ql, qr, t);
        }
        if(qr > mid) {
            mod(u * 2 + 1, mid + 1, r, ql, qr, t);
        }
\end{codeline}
\vspace{-0.9\baselineskip}
\begin{codeline}[firstnumber=47]
        push_up(u);
    }

    Value query(int u, int l, int r, int ql, int qr) {
        if(ql <= l && r <= qr) {
            return node[u].data;
        }
        int mid = l + (r - l) / 2;
        push_down(u, l, r);
        Value res{};
        if(ql <= mid) {
            res = res + query(u * 2, l, mid, ql, qr);
        }
        if(qr > mid) {
            res = res + query(u * 2 + 1, mid + 1, r, ql, qr);
        }
        return res;
    }
    /*
     * 建立[1, n]，所有位置初始化为Value{}
     */
    void build(int n) {
        this->n = n;
        node.assign(4 * n + 4, Node());
        if(n) {
            build(1, 1, n);
        }
    }
    /*
     * 对闭区间[l,r]施加标记t
     */
    void mod(int l, int r, const Tag& t) {
        mod(1, 1, n, l, r, t);
    }
    /*
     * 合并闭区间[l,r]的信息
     */
    Value query(int l, int r) {
        return query(1, 1, n, l, r);
    }
};
\end{codeline}
\begin{warnbox}
Data 的 \texttt{operator+} 与 Tag 的 \texttt{operator+}/\texttt{apply} 必须满足题目对应的幺元与结合律。
\end{warnbox}

\section{动态开点线段树}
\entrymeta{$O(\log V)$}
未开点区间的初值由 \texttt{make\_value(l,\ r)} 给出，默认是 \texttt{Value\{\}}；需要与区间相关的初值时（例如让每个点带上自己的下标），继承并重写它。


\begin{codeline}
class Data {
public:
    int64 sum = 0;
    Data operator+(const Data& other) const {
        return {sum + other.sum};
    }
};

class Tag {
public:
    int64 add = 0;
    Tag operator+(const Tag& other) const {
        return {add + other.add};
    }
    Data apply(const Data& d, int l, int r) const {
        return {d.sum + add * (r - l + 1)};
    }
};

template <typename Value, typename Tag>
class SegmentTree {
public:
    class Node {
    public:
        int ls = -1, rs = -1;
        Value val{};
        Tag tag{};
    };

    vector<Node> node;
    int ncnt = 0;
    int root;
    int le, re;

    /*
     * 未开点区间[l, r]的初值；默认全零
     */
    virtual Value make_value(int, int) {
        return Value{};
    }

    int new_node(int l, int r) {
        int idx = ncnt++;
        node[idx].val = make_value(l, r);
        return idx;
    }
\end{codeline}
\vspace{-0.9\baselineskip}
\begin{codeline}[firstnumber=47]

    void apply(int idx, const Tag& t, int l, int r) {
        node[idx].val = t.apply(node[idx].val, l, r);
        node[idx].tag = node[idx].tag + t;
    }

    void push_down(int l, int r, int u) {
        int mid = l + (r - l) / 2;
        apply(node[u].ls, node[u].tag, l, mid);
        apply(node[u].rs, node[u].tag, mid + 1, r);
        node[u].tag = Tag();
    }

    void mod(int l, int r, int u, const Tag& t, int ql, int qr) {
        if(ql <= l && r <= qr) {
            apply(u, t, l, r);
            return;
        }
        int mid = l + (r - l) / 2;
        if(node[u].ls == -1) {
            node[u].ls = new_node(l, mid);
            node[u].rs = new_node(mid + 1, r);
        }
        push_down(l, r, u);
        if(ql <= mid) {
            mod(l, mid, node[u].ls, t, ql, qr);
        }
        if(qr > mid) {
            mod(mid + 1, r, node[u].rs, t, ql, qr);
        }
        node[u].val = node[node[u].ls].val + node[node[u].rs].val;
    }

    Value query(int l, int r, int u, int ql, int qr) {
        if(u == -1) {
            return make_value(max(l, ql), min(r, qr));
        }
        if(ql <= l && r <= qr) {
            return node[u].val;
        }
        int mid = l + (r - l) / 2;
        Value res{};
        if(ql <= mid) {
            res = res + query(l, mid, node[u].ls, ql, qr);
        }
        if(qr > mid) {
\end{codeline}
\vspace{-0.9\baselineskip}
\begin{codeline}[firstnumber=93]
            res = res + query(mid + 1, r, node[u].rs, ql, qr);
        }
        return node[u].tag.apply(res, max(l, ql), min(r, qr));
    }

    /*
     * 初始化闭区间[l,r]；capacity 按 n * 2 * log V 预留
     */
    void init(int l, int r, int capacity) {
        le = l;
        re = r;
        node.assign(capacity, Node());
        ncnt = 0;
        root = new_node(l, r);
    }

    /*
     * 对闭区间[l,r]施加标记t，按需开点
     */
    void mod(int l, int r, const Tag& t) {
        mod(le, re, root, t, l, r);
    }

    /*
     * 查询闭区间[l,r]的信息，查询不创建节点
     */
    Value query(int l, int r) {
        return query(le, re, root, l, r);
    }
};
\end{codeline}
\begin{warnbox}
\texttt{init} 的 capacity 按 n * 2 * log V 预留，节点池用 \texttt{ncnt} 分配，不 \texttt{push\_back}；缺失子节点按空信息处理。
\end{warnbox}

\section{标记持久化线段树}
\entrymeta{$O(\log n)$}
\begin{codeline}
const int INVALID = -2147483648;

template <class Value, class Tag>
class SegmentTree {
public:
    class Node {
    public:
        int ls = INVALID, rs = INVALID;
        Value val{};
        Tag tag{};
    };

    vector<Node> node;
    int root = INVALID;
    int le = 0, re = 0;

    int new_node() {
        int idx = int(node.size());
        node.emplace_back();
        return idx;
    }

    /*
     * 用本节点标记和左右子节点信息重算本节点信息
     */
    void push_up(int u, int l, int r) {
        Value child{};
        if(l != r && node[u].ls != INVALID) {
            child = node[node[u].ls].val + node[node[u].rs].val;
        }
        node[u].val = node[u].tag.apply(child, l, r);
    }

    void mod(int u, int l, int r, int ql, int qr, const Tag& t) {
        if(ql <= l && r <= qr) {
            node[u].tag = node[u].tag + t;
            push_up(u, l, r);
            return;
        }
        int mid = l + (r - l) / 2;
        if(node[u].ls == INVALID) {
            node[u].ls = new_node();
            node[u].rs = new_node();
        }
        if(ql <= mid) {
            mod(node[u].ls, l, mid, ql, qr, t);
\end{codeline}
\vspace{-0.9\baselineskip}
\begin{codeline}[firstnumber=47]
        }
        if(qr > mid) {
            mod(node[u].rs, mid + 1, r, ql, qr, t);
        }
        push_up(u, l, r);
    }

    Value query(int u, int l, int r, int ql, int qr) {
        if(u == INVALID) {
            return Value();
        }
        if(ql <= l && r <= qr) {
            return node[u].val;
        }
        int mid = l + (r - l) / 2;
        Value res{};
        if(ql <= mid) {
            res = res + query(node[u].ls, l, mid, ql, qr);
        }
        if(qr > mid) {
            res = res + query(node[u].rs, mid + 1, r, ql, qr);
        }
        return node[u].tag.apply(res, max(l, ql), min(r, qr));
    }

    /*
     * 初始化闭区间[l,r]；ncnt为节点池预留量
     */
    void init(int l, int r, int ncnt = 1) {
        le = l;
        re = r;
        node.clear();
        node.reserve(ncnt);
        root = new_node();
    }

    /*
     * 对闭区间[l,r]施加标记t，标记不下传
     */
    void mod(int l, int r, const Tag& t) {
        mod(root, le, re, l, r, t);
    }

    /*
     * 查询闭区间[l,r]的信息
     */
\end{codeline}
\vspace{-0.9\baselineskip}
\begin{codeline}[firstnumber=93]
    Value query(int l, int r) {
        return query(root, le, re, l, r);
    }
};
\end{codeline}
\begin{warnbox}
标记不下传：\texttt{push\_up} 时把本节点累积标记作用在左右儿子合并值上；这是标记持久化，不是懒惰标记。
\end{warnbox}

\section{吉司机线段树}
\entrymeta{$O(\log^2 n)$ 均摊}
\subsection{P6242 【模板】线段树 3（区间最值操作、区间历史最值）}\label{p6242-ux6a21ux677fux7ebfux6bb5ux6811-3ux533aux95f4ux6700ux503cux64cdux4f5cux533aux95f4ux5386ux53f2ux6700ux503c}


\begin{codeline}
const int INVALID = INT_MIN;
const int64 INF = 0x3f3f3f3f3f3f3f3f;
const int N = 500005;
int n, w;

class Tag {
public:
    int64 add, hadd;
    int64 seadd, hseadd;

    Tag operator+(const Tag& other) const {
        Tag res = {
            add + other.add,
            max(hadd, add + other.hadd),
            seadd + other.seadd,
            max(hseadd, seadd + other.hseadd),
        };

        return res;
    }
};

class Data {
public:
    int64 maxi, semax;
    int64 hmaxi;
    int64 span, sum, mcnt;

    Data operator+(const Data& other) const {
        Data res = {
            max(maxi, other.maxi),
            INVALID,
            max(hmaxi, other.hmaxi),
            span + other.span,
            sum + other.sum,
            maxi > other.maxi ? mcnt : other.mcnt,
        };

        if(maxi == other.maxi) {
            res.semax = max(semax, other.semax);
            res.mcnt = mcnt + other.mcnt;
        } else {
            res.semax = max(min(maxi, other.maxi), max(semax, other.semax));
        }

        return res;
\end{codeline}
\vspace{-0.9\baselineskip}
\begin{codeline}[firstnumber=47]
    }

    Data operator+(const Tag& tag) const {
        Data res = {
            maxi + tag.add,
            semax + tag.seadd,
            max(hmaxi, maxi + tag.hadd),
            span,
            sum + tag.add * mcnt + tag.seadd * (span - mcnt),
            mcnt,
        };
        return res;
    }

    void read(int idx) {
        int num;
        cin >> num;
        maxi = hmaxi = sum = num;
        semax = -INF;
        span = mcnt = 1;
    }
};

class Node {
public:
    int ls, rs; // Left Son, Right Son
    Data data;
    Tag tag;
};

class Condition {
public:
};

const Tag TAG_EP = {0, 0, 0, 0};

class SegmentTree {
private:
    static const int MAX_N = N;
    int size, le, re, opl, opr, root;
    Node node[MAX_N * 2];

    void build(int l, int r, int cur) {
        node[cur].tag = TAG_EP;
        if(l == r) {
            node[cur].ls = node[cur].rs = INVALID;
\end{codeline}
\vspace{-0.9\baselineskip}
\begin{codeline}[firstnumber=93]
            node[cur].data.read(l);
            return;
        } else {
            int mid = (l + r) >> 1;
            node[cur].ls = cur + 1;
            node[cur].rs = cur + 2 * (mid - l + 1);
            build(l, mid, node[cur].ls);
            build(mid + 1, r, node[cur].rs);
            push_up(cur);
        }
    }

    void update(int cur, Tag tag) {
        node[cur].tag = node[cur].tag + tag;
        node[cur].data = node[cur].data + tag;
    }

    void push_down(int cur) {
        if(node[node[cur].ls].data.maxi == node[node[cur].rs].data.maxi) {
            update(node[cur].ls, node[cur].tag);
            update(node[cur].rs, node[cur].tag);
        } else if(node[node[cur].ls].data.maxi > node[node[cur].rs].data.maxi) {
            Tag se = node[cur].tag;
            se.add = se.seadd;
            se.hadd = se.hseadd;
            update(node[cur].ls, node[cur].tag);
            update(node[cur].rs, se);
        } else {
            Tag se = node[cur].tag;
            se.add = se.seadd;
            se.hadd = se.hseadd;
            update(node[cur].ls, se);
            update(node[cur].rs, node[cur].tag);
        }
        node[cur].tag = TAG_EP;
    }

    void push_up(int cur) {
        node[cur].data = node[node[cur].ls].data + node[node[cur].rs].data;
    }

    void mod(int l, int r, int cur, const Tag& x) {
        if(opl <= l && r <= opr) {
            update(cur, x);
            return;
        }
\end{codeline}
\vspace{-0.9\baselineskip}
\begin{codeline}[firstnumber=139]
        int mid = (l + r) >> 1;
        push_down(cur);
        if(opl <= l && mid <= opr) {
            update(node[cur].ls, x);
        } else if(max(opl, l) <= min(opr, mid)) {
            mod(l, mid, node[cur].ls, x);
        }
        if(opl <= mid + 1 && r <= opr) {
            update(node[cur].rs, x);
        } else if(max(opl, mid + 1) <= min(opr, r)) {
            mod(mid + 1, r, node[cur].rs, x);
        }
        push_up(cur);
    }

    bool cupdate(int cur, int64 mini) {
        if(node[cur].data.maxi <= mini) {
            return true;
        } else if(node[cur].data.semax < mini) {
            int64 g = max(0ll, node[cur].data.maxi - mini);
            node[cur].data.maxi -= g;
            node[cur].data.sum -= g * node[cur].data.mcnt;
            node[cur].tag.add -= g;
            return true;
        } else {
            return false;
        }
    }

    /*
     * 对闭区间[l,r]执行取min操作
     */
    void cmod(int l, int r, int cur, int64 mini) {
        if(opl <= l && r <= opr && cupdate(cur, mini)) {
            return;
        }
        int mid = (l + r) >> 1;
        push_down(cur);
        if(opl <= l && mid <= opr && cupdate(node[cur].ls, mini)) {
        } else if(max(opl, l) <= min(opr, mid)) {
            cmod(l, mid, node[cur].ls, mini);
        }
        if(opl <= mid + 1 && r <= opr && cupdate(node[cur].rs, mini)) {
        } else if(max(opl, mid + 1) <= min(opr, r)) {
            cmod(mid + 1, r, node[cur].rs, mini);
        }
\end{codeline}
\vspace{-0.9\baselineskip}
\begin{codeline}[firstnumber=185]
        push_up(cur);
    }

    Data query(int l, int r, int cur) {
        if(opl <= l && r <= opr) {
            return node[cur].data;
        }
        int mid = (l + r) >> 1;
        push_down(cur);
        Data resl, resr;
        int state = 3;
        if(opl <= l && mid <= opr) {
            resl = node[node[cur].ls].data;
        } else if(max(opl, l) <= min(opr, mid)) {
            resl = query(l, mid, node[cur].ls);
        } else {
            state ^= 2;
        }
        if(opl <= mid + 1 && r <= opr) {
            resr = node[node[cur].rs].data;
        } else if(max(opl, mid + 1) <= min(opr, r)) {
            resr = query(mid + 1, r, node[cur].rs);
        } else {
            state ^= 1;
        }

        switch(state) {
        case(3):
            return (resl + resr);
        case(2):
            return resl;
        case(1):
            return resr;
        }
        return {};
    }

public:
    /*
     * 建立闭区间[l,r]的线段树
     */
    void build(int l, int r) {
        root = 0;
        le = l;
        re = r;
        size = r - l + 1;
\end{codeline}
\vspace{-0.9\baselineskip}
\begin{codeline}[firstnumber=231]
        build(l, r, root);
    }

    /*
     * 对闭区间[l,r]施加加法标记x
     */
    void mod(int l, int r, const Tag& x) {
        opl = l;
        opr = r;
        mod(le, re, root, x);
    }

    /*
     * 对闭区间[l,r]执行取min操作
     */
    void cmod(int l, int r, int64 mini) {
        opl = l;
        opr = r;
        cmod(le, re, root, mini);
    }

    /*
     * 查询闭区间[l,r]的维护信息
     */
    Data query(int l, int r) {
        opl = l;
        opr = r;
        return query(le, re, root);
    }
};

SegmentTree segmentTree;
\end{codeline}


\subsection{U180387 CTSN loves segment tree}\label{u180387-ctsn-loves-segment-tree}


\begin{codeline}
const int INVALID = INT_MIN;
const int64 INF = 0x3f3f3f3f3f3f3f3f;
const int N = 300005;
int n, w;
int a[N], b[N];

class Tag {
public:
    int64 Aadd;
    int64 seAadd;
    int64 Badd;
    int64 seBadd;

    Tag operator+(const Tag& other) const {
        Tag res = {
            Aadd + other.Aadd,
            seAadd + other.seAadd,
            Badd + other.Badd,
            seBadd + other.seBadd,
        };

        return res;
    }
};

class CTag {
public:
    int64 miniA;
    int64 miniB;
};

class Data {
public:
    int64 Amaxi, seAmax;
    int64 Bmaxi, seBmax;
    int64 c[4];
    // 11: Amaxi Bmaxi

    Data operator+(const Data& other) const {
        Data res = {
            max(Amaxi, other.Amaxi),
            Amaxi == other.Amaxi ? max(seAmax, other.seAmax) : max(min(Amaxi, other.Amaxi), max(seAmax, other.seAmax)),
            max(Bmaxi, other.Bmaxi),
            Bmaxi == other.Bmaxi ? max(seBmax, other.seBmax) : max(min(Bmaxi, other.Bmaxi), max(seBmax, other.seBmax)),
            -INF,
            -INF,
\end{codeline}
\vspace{-0.9\baselineskip}
\begin{codeline}[firstnumber=47]
            -INF,
            -INF};

        if(Amaxi == res.Amaxi) {
            if(Bmaxi == res.Bmaxi) {
                for(int i = 0; i < 4; i++) {
                    res.c[i] = max(res.c[i], c[i]);
                }
            } else {
                res.c[2] = max(res.c[2], c[3]);
                res.c[2] = max(res.c[2], c[2]);
                res.c[0] = max(res.c[0], c[1]);
                res.c[0] = max(res.c[0], c[0]);
            }
        } else {
            if(Bmaxi == res.Bmaxi) {
                res.c[1] = max(res.c[1], c[3]);
                res.c[1] = max(res.c[1], c[1]);
                res.c[0] = max(res.c[0], c[2]);
                res.c[0] = max(res.c[0], c[0]);
            } else {
                for(int i = 0; i < 4; i++) {
                    res.c[0] = max(res.c[0], c[i]);
                }
            }
        }

        if(other.Amaxi == res.Amaxi) {
            if(other.Bmaxi == res.Bmaxi) {
                for(int i = 0; i < 4; i++) {
                    res.c[i] = max(res.c[i], other.c[i]);
                }
            } else {
                res.c[2] = max(res.c[2], other.c[3]);
                res.c[2] = max(res.c[2], other.c[2]);
                res.c[0] = max(res.c[0], other.c[1]);
                res.c[0] = max(res.c[0], other.c[0]);
            }
        } else {
            if(other.Bmaxi == res.Bmaxi) {
                res.c[1] = max(res.c[1], other.c[3]);
                res.c[1] = max(res.c[1], other.c[1]);
                res.c[0] = max(res.c[0], other.c[2]);
                res.c[0] = max(res.c[0], other.c[0]);
            } else {
                for(int i = 0; i < 4; i++) {
\end{codeline}
\vspace{-0.9\baselineskip}
\begin{codeline}[firstnumber=93]
                    res.c[0] = max(res.c[0], other.c[i]);
                }
            }
        }

        return res;
    }

    Data operator+(const Tag& tag) const {
        Data res = {
            Amaxi + tag.Aadd,
            seAmax + tag.seAadd,
            Bmaxi + tag.Badd,
            seBmax + tag.seBadd,
            c[0] == -INF ? -INF : c[0] + tag.seAadd + tag.seBadd,
            c[1] == -INF ? -INF : c[1] + tag.seAadd + tag.Badd,
            c[2] == -INF ? -INF : c[2] + tag.Aadd + tag.seBadd,
            c[3] == -INF ? -INF : c[3] + tag.Aadd + tag.Badd};
        return res;
    }

    void read(int idx) {
        Amaxi = a[idx];
        Bmaxi = b[idx];
        seAmax = seBmax = -INF;
        c[3] = Amaxi + Bmaxi;
        c[0] = c[1] = c[2] = -INF;
    }
};

class Node {
public:
    int ls, rs; // Left Son, Right Son
    Data data;
    Tag tag;
};

class Condition {
public:
};

const Tag TAG_EP = {0, 0, 0, 0};

class SegmentTree {
private:
    static const int MAX_N = N;
\end{codeline}
\vspace{-0.9\baselineskip}
\begin{codeline}[firstnumber=139]
    int size, le, re, opl, opr, root;
    /*
     * le, re: Left Edge, Right Edge
     * opl, opr: Operation l, r
     * root: default 0
     */
    Node node[MAX_N * 2];

    void build(int l, int r, int cur) {
        node[cur].tag = TAG_EP;
        if(l == r) {
            node[cur].ls = node[cur].rs = INVALID;
            node[cur].data.read(l);
            return;
        } else {
            int mid = (l + r) >> 1;
            node[cur].ls = cur + 1;
            node[cur].rs = cur + 2 * (mid - l + 1);
            build(l, mid, node[cur].ls);
            build(mid + 1, r, node[cur].rs);
            push_up(cur);
        }
    }

    void update(int cur, Tag tag) {
        node[cur].tag = node[cur].tag + tag;
        node[cur].data = node[cur].data + tag;
    }

    void push_down(int cur) {
        if(node[node[cur].ls].data.Amaxi == node[node[cur].rs].data.Amaxi && node[node[cur].ls].data.Bmaxi == node[node[cur].rs].data.Bmaxi) {
            update(node[cur].ls, node[cur].tag);
            update(node[cur].rs, node[cur].tag);
        } else {
            Tag lt = node[cur].tag, rt = node[cur].tag;
            if(node[node[cur].ls].data.Amaxi != node[node[cur].rs].data.Amaxi) {
                if(node[node[cur].ls].data.Amaxi < node[node[cur].rs].data.Amaxi) {
                    lt.Aadd = lt.seAadd;
                } else {
                    rt.Aadd = rt.seAadd;
                }
            }

            if(node[node[cur].ls].data.Bmaxi != node[node[cur].rs].data.Bmaxi) {
                if(node[node[cur].ls].data.Bmaxi < node[node[cur].rs].data.Bmaxi) {
                    lt.Badd = lt.seBadd;
\end{codeline}
\vspace{-0.9\baselineskip}
\begin{codeline}[firstnumber=185]
                } else {
                    rt.Badd = rt.seBadd;
                }
            }

            update(node[cur].ls, lt);
            update(node[cur].rs, rt);
        }
        node[cur].tag = TAG_EP;
    }

    void push_up(int cur) {
        node[cur].data = node[node[cur].ls].data + node[node[cur].rs].data;
    }

    void mod(int l, int r, int cur, const Tag& x) {
        if(opl <= l && r <= opr) {
            update(cur, x);
            return;
        }
        int mid = (l + r) >> 1;
        push_down(cur);
        if(opl <= l && mid <= opr) {
            update(node[cur].ls, x);
        } else if(max(opl, l) <= min(opr, mid)) {
            mod(l, mid, node[cur].ls, x);
        }
        if(opl <= mid + 1 && r <= opr) {
            update(node[cur].rs, x);
        } else if(max(opl, mid + 1) <= min(opr, r)) {
            mod(mid + 1, r, node[cur].rs, x);
        }
        push_up(cur);
    }

    bool cupdate(int cur, const CTag& mini) {
        if(node[cur].data.Amaxi <= mini.miniA && node[cur].data.Bmaxi <= mini.miniB) {
            return true;
        } else if(node[cur].data.seAmax < mini.miniA && node[cur].data.seBmax < mini.miniB) {
            int64 gA = max(0ll, node[cur].data.Amaxi - mini.miniA);
            int64 gB = max(0ll, node[cur].data.Bmaxi - mini.miniB);
            node[cur].data.Amaxi -= gA;
            node[cur].tag.Aadd -= gA;
            if(node[cur].data.c[3] != -INF) {
                node[cur].data.c[3] -= gA;
            }
\end{codeline}
\vspace{-0.9\baselineskip}
\begin{codeline}[firstnumber=231]
            if(node[cur].data.c[2] != -INF) {
                node[cur].data.c[2] -= gA;
            }

            node[cur].data.Bmaxi -= gB;
            node[cur].tag.Badd -= gB;
            if(node[cur].data.c[3] != -INF) {
                node[cur].data.c[3] -= gB;
            }
            if(node[cur].data.c[1] != -INF) {
                node[cur].data.c[1] -= gB;
            }
            return true;
        } else {
            return false;
        }
    }

    /*
     * 对闭区间[l,r]执行取min操作
     */
    void cmod(int l, int r, int cur, const CTag& mini) {
        if(opl <= l && r <= opr && cupdate(cur, mini)) {
            return;
        }
        int mid = (l + r) >> 1;
        push_down(cur);
        if(opl <= l && mid <= opr && cupdate(node[cur].ls, mini)) {
        } else if(max(opl, l) <= min(opr, mid)) {
            cmod(l, mid, node[cur].ls, mini);
        }
        if(opl <= mid + 1 && r <= opr && cupdate(node[cur].rs, mini)) {
        } else if(max(opl, mid + 1) <= min(opr, r)) {
            cmod(mid + 1, r, node[cur].rs, mini);
        }
        push_up(cur);
    }

    Data query(int l, int r, int cur) {
        if(opl <= l && r <= opr) {
            return node[cur].data;
        }
        int mid = (l + r) >> 1;
        push_down(cur);
        Data resl, resr;
        int state = 3;
\end{codeline}
\vspace{-0.9\baselineskip}
\begin{codeline}[firstnumber=277]
        if(opl <= l && mid <= opr) {
            resl = node[node[cur].ls].data;
        } else if(max(opl, l) <= min(opr, mid)) {
            resl = query(l, mid, node[cur].ls);
        } else {
            state ^= 2;
        }
        if(opl <= mid + 1 && r <= opr) {
            resr = node[node[cur].rs].data;
        } else if(max(opl, mid + 1) <= min(opr, r)) {
            resr = query(mid + 1, r, node[cur].rs);
        } else {
            state ^= 1;
        }

        switch(state) {
        case(3):
            return (resl + resr);
        case(2):
            return resl;
        case(1):
            return resr;
        }
        return {};
    }

public:
    /*
     * 建立闭区间[l,r]的线段树
     */
    void build(int l, int r) {
        /*
         * 极端情况下，线段树可能只有一个节点
         * 在这种情况下上述写法会报错
         */
        root = 0;
        le = l;
        re = r;
        size = r - l + 1;
        build(l, r, root);
    }

    /*
     * 对闭区间[l,r]施加加法标记x
     */
    void mod(int l, int r, const Tag& x) {
\end{codeline}
\vspace{-0.9\baselineskip}
\begin{codeline}[firstnumber=323]
        opl = l;
        opr = r;
        mod(le, re, root, x);
    }

    /*
     * 对闭区间[l,r]执行取min操作
     */
    void cmod(int l, int r, const CTag& mini) {
        opl = l;
        opr = r;
        cmod(le, re, root, mini);
    }

    /*
     * 查询闭区间[l,r]的维护信息
     */
    Data query(int l, int r) {
        opl = l;
        opr = r;
        return query(le, re, root);
    }
};

SegmentTree segmentTree;
\end{codeline}
\begin{warnbox}
吉司机线段树：\texttt{push\_down} 按左右最大值分配全量/非最大值标记；\texttt{semax} 需要 \texttt{max\_second \textless{} x \textless{} max\_first}。
\end{warnbox}

\section{李超线段树}
\entrymeta{$O(\log^2 n)$}
\subsection{P4097 【模板】李超线段树}\label{p4097-ux6a21ux677fux674eux8d85ux7ebfux6bb5ux6811}


\begin{codeline}
const double FINF = 1e100;
const double EPS = 1e-9;
const int MODER1 = 39989;
const int MODER2 = 1000000000;

int n;
int last_ans = 0;

class Line {
public:
    double k, b;
    int idx;

    inline double operator()(double x) const {
        return k * x + b;
    }
};

class LiChaoSegmentTreeNode {
public:
    int range_l, range_r;
    int l, r;
    Line g;
};

class LiChaoSegmentTree {
private:
    static const int NO_TAG = 0xffffffff;
    static const int NO_SON = -1;
    static const int MAX_N = 40001;
    int size;
    int root = 0;
    int node_used = 1; // 0 for root
    int op_range_l, op_range_r;
    int ans_idx;
    double ans;
    LiChaoSegmentTreeNode node[MAX_N * 2];

    int getNode() {
        return node_used++;
    }

    void build(int range_l, int range_r, int cur) {
        node[cur].range_l = range_l;
        node[cur].range_r = range_r;
        node[cur].l = NO_SON;
\end{codeline}
\vspace{-0.9\baselineskip}
\begin{codeline}[firstnumber=47]
        node[cur].r = NO_SON;
        node[cur].g = {FINF, -FINF, -1};
        if(range_r == range_l) {
            return;
        } else {
            int mid = (range_r + range_l) >> 1;
            node[cur].l = getNode();
            node[cur].r = getNode();
            build(range_l, mid, node[cur].l);
            build(mid + 1, range_r, node[cur].r);
        }
    }

    void update(int cur, Line f) {
        if(cur == NO_SON) {
            return;
        }
        if(node[cur].g.idx == -1) {
            node[cur].g = f;
        } else {

            int mid = (node[cur].range_r + node[cur].range_l) >> 1;
            double mid_f, mid_g, l_f, l_g, r_f, r_g;
            mid_f = f(mid);
            mid_g = node[cur].g(mid);

            if(mid_f - mid_g > EPS || abs(mid_f - mid_g) < EPS && node[cur].g.idx > f.idx) {
                swap(node[cur].g, f);
            }

            if(node[cur].range_l == node[cur].range_r) {
                return;
            }
            if(f.k == node[cur].g.k && f.b == node[cur].g.b) {
                return;
            }
            l_f = f(node[cur].range_l);
            r_f = f(node[cur].range_r);
            l_g = node[cur].g(node[cur].range_l);
            r_g = node[cur].g(node[cur].range_r);
            if(r_f - r_g > EPS || abs(r_f - r_g) < EPS) {
                update(node[cur].r, f);
            }
            if(l_f - l_g > EPS || abs(l_f - l_g) < EPS) {
                update(node[cur].l, f);
            }
\end{codeline}
\vspace{-0.9\baselineskip}
\begin{codeline}[firstnumber=93]
        }
    }

    void add_line(int cur, Line f) {
        if(cur == NO_SON) {
            return;
        }
        if(max(op_range_l, node[cur].range_l) <= min(op_range_r, node[cur].range_r)) {
            if(op_range_l <= node[cur].range_l && node[cur].range_r <= op_range_r) {
                update(cur, f);
            } else {
                add_line(node[cur].l, f);
                add_line(node[cur].r, f);
            }
        }
    }

    void query(int cur, int x) {
        if(cur == NO_SON) {
            return;
        }
        if(node[cur].range_l <= x && x <= node[cur].range_r) {
            if(node[cur].g.idx != -1) {
                double y_x;
                y_x = node[cur].g(x);
                if(y_x - ans > EPS || abs(y_x - ans) < EPS && ans_idx > node[cur].g.idx) {
                    ans = y_x;
                    ans_idx = node[cur].g.idx;
                }
            }
            query(node[cur].l, x);
            query(node[cur].r, x);
        }
    }

public:
    /*
     * 初始化横坐标闭区间[range_l,range_r]
     */
    void build(int range_l, int range_r) {
        root = 0;
        node_used = 1;
        build(range_l, range_r, root);
        size = range_r - range_l + 1;
    }

\end{codeline}
\vspace{-0.9\baselineskip}
\begin{codeline}[firstnumber=139]
    /*
     * 插入端点为(x1,y1)、(x2,y2)的线段，编号为idx
     */
    void add_line(int x1, int y1, int x2, int y2, int idx) {
        op_range_l = min(x1, x2);
        op_range_r = max(x1, x2);
        Line f;
        if(x1 == x2) {
            f = {0., (double)(max(y2, y1)), idx};
        } else {
            double k = ((double)(y2 - y1)) / (x2 - x1);
            f = {k, y1 - k * (x1), idx};
        }
        add_line(root, f);
    }

    /*
     * 返回横坐标x处最高线段的编号，高度相同取编号较小者
     */
    int query(int x) {
        ans = -FINF;
        ans_idx = -1;
        query(root, x);
        if(ans_idx == -1) {
            return 0;
        }
        return ans_idx;
    }
};

LiChaoSegmentTree segmentTree;
\end{codeline}


\subsection{离散坐标：区间最小值}\label{ux79bbux6563ux5750ux6807ux533aux95f4ux6700ux5c0fux503c}

用于在线段覆盖的离散坐标上查询区间最小值，与上面的单点最大线段编号版本独立使用。 \texttt{a{[}1..n{]}} 非递减，允许重复坐标；操作参数是下标，满足 \texttt{1\ \textless{}=\ l\ \textless{}=\ r\ \textless{}=\ n}。 初始水平线也参与取最小值；\texttt{k*x+b} 的结果需在 int64 范围内，中间乘法使用 int128。 建树 O(n)，区间加线段 O(log² n)，区间查询 O(log n)，空间 O(n)。


\begin{codeline}
class Line {
public:
    int64 k = 0, b = 0;
    int64 operator()(int64 x) const {
        return int64(int128(k) * x + b);
    }
};

class LiChaoSegmentTree {
public:
    class Node {
    public:
        Line g;
        int64 minimum;
    };
    int n = 0;
    vector<int64> x;
    vector<Node> node;

    void push_up(int u, int l, int r) {
        node[u].minimum = min(node[u].g(x[l]), node[u].g(x[r]));
        if(l < r) {
            node[u].minimum = min(node[u].minimum, min(node[u * 2].minimum, node[u * 2 + 1].minimum));
        }
    }
    void update(int u, int l, int r, Line f) {
        int mid = (l + r) / 2;
        if(f(x[mid]) < node[u].g(x[mid])) {
            swap(f, node[u].g);
        }
        if(l < r) {
            if(f(x[l]) < node[u].g(x[l])) {
                update(u * 2, l, mid, f);
            } else if(f(x[r]) < node[u].g(x[r])) {
                update(u * 2 + 1, mid + 1, r, f);
            }
        }
        push_up(u, l, r);
    }
    void add_line(int u, int l, int r, int ql, int qr, Line f) {
        if(ql <= l && r <= qr) {
            update(u, l, r, f);
            return;
        }
        int mid = (l + r) / 2;
        if(ql <= mid) {
\end{codeline}
\vspace{-0.9\baselineskip}
\begin{codeline}[firstnumber=47]
            add_line(u * 2, l, mid, ql, qr, f);
        }
        if(qr > mid) {
            add_line(u * 2 + 1, mid + 1, r, ql, qr, f);
        }
        push_up(u, l, r);
    }
    int64 range_query(int u, int l, int r, int ql, int qr) const {
        if(ql <= l && r <= qr) {
            return node[u].minimum;
        }
        int64 res = min(node[u].g(x[max(l, ql)]), node[u].g(x[min(r, qr)]));
        int mid = (l + r) / 2;
        if(ql <= mid) {
            res = min(res, range_query(u * 2, l, mid, ql, qr));
        }
        if(qr > mid) {
            res = min(res, range_query(u * 2 + 1, mid + 1, r, ql, qr));
        }
        return res;
    }
    /*
     * 由非递减横坐标a[1..n]建树，以水平线value初始化
     */
    void build(const vector<int64>& a, int64 value) {
        x = a;
        n = int(a.size()) - 1;
        node.assign(4 * n + 4, Node{{0, value}, value});
    }
    /*
     * 在下标闭区间[l,r]加入直线kx+b
     */
    void add_line(int l, int r, int64 k, int64 b) {
        add_line(1, 1, n, l, r, {k, b});
    }
    /*
     * 查询下标闭区间[l,r]内的最小值
     */
    int64 range_query(int l, int r) const {
        return range_query(1, 1, n, l, r);
    }
};
\end{codeline}
\begin{warnbox}
用 double 比较，\texttt{EPS} 过小时极端坐标可能选错线段；需要精确时改成 int128 交叉相乘。
\end{warnbox}

\section{线段树合并}
\entrymeta{$O(n\log n)$}
本例只做\textbf{单点加、整树查询、破坏性合并}；mod 的 l、r 必须相等，不能用于区间加。 合并后子树节点可能被父树复用，子树答案在 dfs3 返回前保存，之后不要再查询已合并的子树。 固定节点池需按 O((n+m) log V) 估算；深链上的 DFS 仍需按赛题栈限制改成迭代。

\subsection{P4556 {[}Vani有约会{]} 雨天的尾巴 /【模板】线段树合并}\label{p4556-vaniux6709ux7ea6ux4f1a-ux96e8ux5929ux7684ux5c3eux5df4-ux6a21ux677fux7ebfux6bb5ux6811ux5408ux5e76}


\begin{codeline}
const int INF = 0x3f3f3f3f;
const int INVALID = -2147483648;
const int N = 100005;
int n, m, moder;

class Tag {
public:
    int add;
    Tag operator+(const Tag& other) {
        return {add + other.add};
    }
};

class Data {
public:
    int maxi, idx;

    Data operator+(const Data& other) {
        if(maxi >= other.maxi) {
            return *this;
        } else {
            return other;
        }
    }

    void read(int i) {
        idx = i;
        cin >> maxi;
    }
};

inline Data merge(int l, int r, const Data& d, const Tag& tag) {
    return {
        d.maxi + tag.add,
        d.idx};
}

class SegmentNode {
public:
    int ls, rs; // Left Son, Right Son
    Data data;
    Tag tag;
};

const Data DATA_EP = {0, 0};
const Tag TAG_EP = {0};
\end{codeline}
\vspace{-0.9\baselineskip}
\begin{codeline}[firstnumber=47]

SegmentNode stNode[200 * N];
int stNode_cnt;

class SegmentTree {
private:
    static const int MAX_N = N;
    int size, le, re, opl, opr, root;

    int getNode(int idx) {
        stNode[stNode_cnt] = {
            INVALID, INVALID, {0, idx}, {0}};
        return stNode_cnt++;
    }

    void push_up(int cur) {
        if(stNode[cur].ls != INVALID && stNode[cur].rs != INVALID) {
            stNode[cur].data = stNode[stNode[cur].ls].data + stNode[stNode[cur].rs].data;
        } else if(stNode[cur].ls != INVALID) {
            stNode[cur].data = stNode[stNode[cur].ls].data;
        } else if(stNode[cur].rs != INVALID) {
            stNode[cur].data = stNode[stNode[cur].rs].data;
        } else {
            return;
        }
    }

    void mod(int l, int r, int cur, const Tag& x) {
        if(opl <= l && r <= opr) {
            stNode[cur].data = merge(l, r, stNode[cur].data, x);
            stNode[cur].tag = stNode[cur].tag + x;
        } else {
            int mid = (l + r) >> 1;
            if(stNode[cur].ls == INVALID) {
                stNode[cur].ls = getNode(l);
                stNode[cur].rs = getNode(mid + 1);
            }
            if(max(opl, l) <= min(opr, mid)) {
                mod(l, mid, stNode[cur].ls, x);
            }
            if(max(opl, mid + 1) <= min(opr, r)) {
                mod(mid + 1, r, stNode[cur].rs, x);
            }
            push_up(cur);
        }
    }
\end{codeline}
\vspace{-0.9\baselineskip}
\begin{codeline}[firstnumber=93]

    Data query(int l, int r, int cur) {
        if(cur == INVALID) {
            return DATA_EP;
        }
        if(opl <= l && r <= opr) {
            return stNode[cur].data;
        } else if(cur != INVALID) {
            int mid = (l + r) >> 1;
            bool dual = max(opl, l) <= min(opr, mid) && max(opl, mid + 1) <= min(opr, r);

            Data res;
            if(dual) {
                res = query(l, mid, stNode[cur].ls) +
                      query(mid + 1, r, stNode[cur].rs);
            } else if(max(opl, l) <= min(opr, mid)) {
                res = query(l, mid, stNode[cur].ls);
            } else {
                res = query(mid + 1, r, stNode[cur].rs);
            }
            return merge(max(opl, l), min(opr, r), res, stNode[cur].tag);
        } else {
            return DATA_EP;
        }
    }

    int dfsMerge(int l, int r, int curS, int curO, const SegmentTree& other) {
        if(curS == INVALID && curO == INVALID) {
            return INVALID;
        } else if(curS != INVALID && curO != INVALID) {
            if(l == r) {
                stNode[curS].data.maxi += stNode[curO].data.maxi;
            } else {
                int mid = (l + r) >> 1;
                stNode[curS].ls = dfsMerge(l, mid, stNode[curS].ls, stNode[curO].ls, other);
                stNode[curS].rs = dfsMerge(mid + 1, r, stNode[curS].rs, stNode[curO].rs, other);
                push_up(curS);
            }
            return curS;
        } else if(curS != INVALID) {
            return curS;
        } else {
            return curO;
        }
    }

\end{codeline}
\vspace{-0.9\baselineskip}
\begin{codeline}[firstnumber=139]
public:
    /*
     * 建立值域闭区间[l,r]的空树
     */
    void build(int l, int r) {
        root = getNode(l);
        le = l;
        re = r;
        size = r - l + 1;
    }

    /*
     * 更新闭区间[l,r]，此例调用时l=r
     */
    void mod(int l, int r, const Tag& x) {
        opl = l;
        opr = r;
        mod(le, re, root, x);
    }

    /*
     * 查询闭区间[l,r]的维护信息
     */
    Data query(int l, int r) {
        opl = l;
        opr = r;
        return query(le, re, root);
    }

    void operator+=(const SegmentTree& other) {
        root = dfsMerge(le, re, root, other.root, other);
    }
};

class HLDTreeNode {
public:
    vector<int> son;
    int dfn, fa, w, link_head, d;
    int heavy, sub_dfn;
    int ans;
    SegmentTree st;
};

class HLDTree {
private:
    static const int MAX_N = N;
\end{codeline}
\vspace{-0.9\baselineskip}
\begin{codeline}[firstnumber=185]
    int dfn = 1;
    HLDTreeNode hldNode[MAX_N];

    void dfs1(int cur, int fa, int d) {
        hldNode[cur].fa = fa;
        hldNode[cur].w = 1;
        hldNode[cur].d = d;
        hldNode[cur].st.build(1, N - 1);
        int heavy = -1, w = 0;
        for(int son : hldNode[cur].son) {
            if(son != fa) {
                dfs1(son, cur, d + 1);
                hldNode[cur].w += hldNode[son].w;
                if(hldNode[son].w > w) {
                    w = hldNode[son].w;
                    heavy = son;
                }
            }
        }
        if(heavy > 0) {
            hldNode[cur].heavy = heavy;
        }
    }

    void dfs2(int cur, int fa, int link_head) {
        inv_dfn[dfn] = cur;
        hldNode[cur].sub_dfn = dfn;
        hldNode[cur].dfn = dfn++;
        hldNode[cur].link_head = link_head;
        if(hldNode[cur].heavy > 0) {
            dfs2(hldNode[cur].heavy, cur, link_head);
            hldNode[cur].sub_dfn = hldNode[hldNode[cur].heavy].sub_dfn;
        }
        for(int son : hldNode[cur].son) {
            if(son != fa && son != hldNode[cur].heavy) {
                dfs2(son, cur, son);
                hldNode[cur].sub_dfn = hldNode[son].sub_dfn;
            }
        }
    }

    int find_lca(int x, int y) {
        while(hldNode[x].link_head != hldNode[y].link_head) { // 先找LCA和计算st之间的距离
            if(hldNode[hldNode[x].link_head].d < hldNode[hldNode[y].link_head].d) {
                swap(x, y);
            }
\end{codeline}
\vspace{-0.9\baselineskip}
\begin{codeline}[firstnumber=231]
            x = hldNode[hldNode[x].link_head].fa;
        }
        return hldNode[x].d < hldNode[y].d ? x : y;
    }

public:
    int root = 1;
    int inv_dfn[MAX_N];

    /*
     * 添加无向边(u,v)
     */
    void add_edge(int u, int v) {
        hldNode[u].son.push_back(v);
        hldNode[v].son.push_back(u);
    }

    /*
     * 以tree_root为根预处理树链信息
     */
    void build(int tree_root) {
        root = tree_root;
        dfs1(root, 0, 1);
        dfs2(root, 0, root);
    }

    /*
     * 在x到y的路径上加入颜色z，使用树上差分
     */
    void path_add(int x, int y, int z) {

        int lca = find_lca(x, y);
        hldNode[x].st.mod(z, z, {1});
        hldNode[y].st.mod(z, z, {1});
        hldNode[lca].st.mod(z, z, {-1});
        if(hldNode[lca].fa) {
            hldNode[hldNode[lca].fa].st.mod(z, z, {-1});
        }
    }

    /*
     * 自底向上合并子树线段树，合并到父亲前保存当前点答案
     */
    void dfs3(int cur) {
        for(int son : hldNode[cur].son) {
            if(son != hldNode[cur].fa) {
\end{codeline}
\vspace{-0.9\baselineskip}
\begin{codeline}[firstnumber=277]
                dfs3(son);
                hldNode[cur].st += hldNode[son].st;
            }
        }
        Data res = hldNode[cur].st.query(1, N - 1);
        hldNode[cur].ans = res.maxi == 0 ? 0 : res.idx;
    }

    int query(int cur) {
        return hldNode[cur].ans;
    }

};

HLDTree hldTree;
\end{codeline}
\begin{warnbox}
节点池按 \texttt{n + 4m log V} 估算；链状树需注意递归深度。
\end{warnbox}

\section{可持久化线段树}
\entrymeta{$O(\log n)$}
\begin{codeline}
class SegmentTree {
public:
    class Node {
    public:
        int ls = 0, rs = 0, cnt = 0;
    };
    vector<Node> node;
    vector<int> root;
    vector<int64> value;

    int mod(int u, int l, int r, int pos) {
        int v = int(node.size());
        node.push_back(node[u]);
        node[v].cnt++;
        if(l != r) {
            int mid = l + (r - l) / 2;
            if(pos <= mid) {
                node[v].ls = mod(node[u].ls, l, mid, pos);
            } else {
                node[v].rs = mod(node[u].rs, mid + 1, r, pos);
            }
        }
        return v;
    }
    /*
     * 将a[1..n]离散化并建立前缀版本
     */
    void build(const vector<int64>& a) {
        int n = int(a.size()) - 1;
        value.assign(a.begin() + 1, a.end());
        sort(value.begin(), value.end());
        value.erase(unique(value.begin(), value.end()), value.end());
        node.assign(1, Node());
        root.assign(n + 1, 0);
        int logn = 0;
        while((1LL << logn) < int(value.size())) {
            logn++;
        }
        node.reserve(1 + size_t(n) * (logn + 1));
        for(int i = 1; i <= n; i++) {
            int pos = lower_bound(value.begin(), value.end(), a[i]) - value.begin();
            root[i] = mod(root[i - 1], 0, int(value.size()) - 1, pos);
        }
    }
    /*
     * 查询闭区间[l,r]中第k小的原值，k从1开始
\end{codeline}
\vspace{-0.9\baselineskip}
\begin{codeline}[firstnumber=47]
     */
    int64 query(int l, int r, int k) const {
        int u = root[r], v = root[l - 1], le = 0, re = int(value.size()) - 1;
        while(le != re) {
            int mid = le + (re - le) / 2;
            int cnt = node[node[u].ls].cnt - node[node[v].ls].cnt;
            if(k <= cnt) {
                u = node[u].ls;
                v = node[v].ls;
                re = mid;
            } else {
                k -= cnt;
                u = node[u].rs;
                v = node[v].rs;
                le = mid + 1;
            }
        }
        return value[le];
    }
};
\end{codeline}


\subsection{P3293 {[}SCOI2016{]} 美味}\label{p3293-scoi2016-ux7f8eux5473}


\begin{codeline}
const int64 MAX_H = 18;
const int N = 200005;
const int M = 100005;
const int64 INVALID = 0;

class Tag {
public:
    int add;
    Tag operator+(const Tag& other) {
        return {add + other.add};
    }
};

class Data {
public:
    int64 sum;
    Data operator+(const Data& other) {
        return {sum + other.sum};
    }
};

inline Data merge(int64 l, int64 r, const Data& d, const Tag& tag) {
    return {d.sum + (r - l + 1) * tag.add};
}

class Node {
public:
    int64 ls, rs;
    Data data;
    Tag tag;
};

const Data DATA_EP = {0};
const Tag TAG_EP = {0};

Node seg_node[N * MAX_H];
int node_cnt = 1;

inline int getNode() {
    return node_cnt++;
}

class SegmentTree {
public:
    static const int64 MAX_N = N;
    int64 size, le, re, opl, opr;
\end{codeline}
\vspace{-0.9\baselineskip}
\begin{codeline}[firstnumber=47]
    Node* node = seg_node;
    int root[N];

    void push_down(int64 l, int64 r, int64 cur) {
        Node& ls = node[node[cur].ls];
        Node& rs = node[node[cur].rs];
        int64 mid = (l + r) >> 1;
        ls.data = merge(l, mid, ls.data, node[cur].tag);
        ls.tag = ls.tag + node[cur].tag;
        rs.data = merge(mid + 1, r, rs.data, node[cur].tag);
        rs.tag = rs.tag + node[cur].tag;
        node[cur].tag = TAG_EP;
    }

    void push_up(int64 cur) {
        node[cur].data = node[node[cur].ls].data + node[node[cur].rs].data;
    }

    void mod(int64 l, int64 r, int64 cur, const int64 lst, const Tag& x) {
        if(opl <= l && r <= opr) {
            node[cur].data = merge(l, r, node[lst].data, x);
            node[cur].tag = node[lst].tag + x;
        } else {
            int64 mid = (l + r) >> 1;
            push_down(l, r, cur);
            if(max(opl, l) <= min(opr, mid)) {
                node[cur].ls = getNode();
                mod(l, mid, node[cur].ls, node[lst].ls, x);
            } else {
                node[cur].ls = node[lst].ls;
            }
            if(max(opl, mid + 1) <= min(opr, r)) {
                node[cur].rs = getNode();
                mod(mid + 1, r, node[cur].rs, node[lst].rs, x);
            } else {
                node[cur].rs = node[lst].rs;
            }
            push_up(cur);
        }
    }

    Data query(int64 l, int64 r, int64 cur, const int lst) {
        if(cur != INVALID) {
            if(opl <= l && r <= opr) {
                return {node[cur].data.sum - node[lst].data.sum};
            } else {
\end{codeline}
\vspace{-0.9\baselineskip}
\begin{codeline}[firstnumber=93]
                int64 mid = (l + r) >> 1;
                bool dual = max(opl, l) <= min(opr, mid) && max(opl, mid + 1) <= min(opr, r);

                Data res;
                if(dual) {
                    res = query(l, mid, node[cur].ls, node[lst].ls) +
                          query(mid + 1, r, node[cur].rs, node[lst].rs);
                } else if(max(opl, l) <= min(opr, mid)) {
                    res = query(l, mid, node[cur].ls, node[lst].ls);
                } else {
                    res = query(mid + 1, r, node[cur].rs, node[lst].rs);
                }
                return merge(max(opl, l), min(opr, r), res, node[cur].tag);
            }
        } else {
            return DATA_EP;
        }
    }

    /*
     * 初始化值域闭区间[l,r]
     */
    void build(int64 l, int64 r) {
        le = l;
        re = r;
        size = r - l + 1;
    }

    /*
     * 基于版本ver修改闭区间[l,r]，写入版本cur
     */
    void mod_from(int cur, int ver, int64 l, int64 r, const Tag& x) {
        opl = l;
        opr = r;
        if(!root[cur]) {
            root[cur] = getNode();
        }
        mod(le, re, root[cur], root[ver], x);
    }

    /*
     * 查询版本cur减去版本ver后的闭区间[l,r]信息
     */
    Data query_from(int cur, int ver, int64 l, int64 r) {
        opl = l;
        opr = r;
\end{codeline}
\vspace{-0.9\baselineskip}
\begin{codeline}[firstnumber=139]
        return query(le, re, root[cur], root[ver]);
    }
};

SegmentTree sgt;

int main() {
    ios::sync_with_stdio(false);
    cin.tie(nullptr);
    int n, m, a;
    cin >> n >> m;
    sgt.build(0, M - 1);
    for(int i = 1; i <= n; i++) {
        cin >> a;
        sgt.mod_from(i, i - 1, a, a, {1});
    }

    int b, x, vl, vr;
    int l, r, sum;
    for(int i = 1; i <= m; i++) {
        cin >> b >> x >> vl >> vr;
        sum = 0;
        for(int j = 1; j <= MAX_H; j++) {
            if((1 << (MAX_H - j)) & b) {
                int raw_l = sum - x;
                int raw_r = sum + (1 << (MAX_H - j)) - x - 1;
                Data res = DATA_EP;
                if(raw_l <= raw_r && raw_r >= 0 && raw_l <= M - 1) {
                    l = max(0, raw_l);
                    r = min(M - 1, raw_r);
                    res = sgt.query_from(vr, vl - 1, l, r);
                }
                if(res.sum > 0) {
                } else {
                    sum += 1 << (MAX_H - j);
                }
            } else {
                int raw_l = sum + (1 << (MAX_H - j)) - x;
                int raw_r = sum + (1 << (MAX_H - j + 1)) - x - 1;
                Data res = DATA_EP;
                if(raw_l <= raw_r && raw_r >= 0 && raw_l <= M - 1) {
                    l = max(0, raw_l);
                    r = min(M - 1, raw_r);
                    res = sgt.query_from(vr, vl - 1, l, r);
                }
                if(res.sum > 0) {
\end{codeline}
\vspace{-0.9\baselineskip}
\begin{codeline}[firstnumber=185]
                    sum += 1 << (MAX_H - j);
                } else {
                }
            }
        }
        cout << (sum ^ b) << '\n';
    }
}
\end{codeline}
\begin{warnbox}
版本编号从 0 开始；区间查询时先判断真实区间与值域是否有交，再钳位。
\end{warnbox}

\section{可持久化数组}
\entrymeta{$O(\log n)$}
\subsection{P3919 【模板】可持久化线段树 1（可持久化数组）}\label{p3919-ux6a21ux677fux53efux6301ux4e45ux5316ux7ebfux6bb5ux6811-1ux53efux6301ux4e45ux5316ux6570ux7ec4}


\begin{codeline}
template <class Data = int>
class PersistentArray {
public:
    class Node {
    public:
        int ls = 0, rs = 0;
        Data data{};
    };
    vector<Node> node;
    vector<int> version;
    int n = 0;

    /*
     * 由a[1..n]建立版本0，capacity仅用于预留节点空间
     */
    int build(int l, int r, const vector<Data>& a) {
        int u = int(node.size());
        node.emplace_back();
        if(l == r) {
            node[u].data = a[l];
        } else {
            int mid = l + (r - l) / 2;
            node[u].ls = build(l, mid, a);
            node[u].rs = build(mid + 1, r, a);
        }
        return u;
    }
    int modify(int u, int l, int r, int pos, const Data& x) {
        int v = int(node.size());
        node.push_back(node[u]);
        if(l == r) {
            node[v].data = x;
        } else {
            int mid = l + (r - l) / 2;
            if(pos <= mid) {
                node[v].ls = modify(node[u].ls, l, mid, pos, x);
            } else {
                node[v].rs = modify(node[u].rs, mid + 1, r, pos, x);
            }
        }
        return v;
    }
    /*
     * 由a[1..n]建立版本0，capacity仅用于预留节点空间
     */
    void build(const vector<Data>& a, size_t capacity = 0) {
\end{codeline}
\vspace{-0.9\baselineskip}
\begin{codeline}[firstnumber=47]
        n = int(a.size()) - 1;
        node.clear();
        version.clear();
        node.reserve(capacity);
        version.push_back(build(1, n, a));
    }
    /*
     * 复制版本ver，返回新版本号
     */
    int copy(int ver) {
        version.push_back(version[ver]);
        return int(version.size()) - 1;
    }
    /*
     * 基于版本ver修改pos处的值，返回新版本号
     */
    int modify(int ver, int pos, const Data& x) {
        int root = modify(version[ver], 1, n, pos, x);
        version.push_back(root);
        return int(version.size()) - 1;
    }
    /*
     * 查询版本ver中pos处的值，不修改历史版本
     */
    Data query(int ver, int pos) const {
        int u = version[ver], l = 1, r = n;
        while(l != r) {
            int mid = l + (r - l) / 2;
            if(pos <= mid) {
                u = node[u].ls;
                r = mid;
            } else {
                u = node[u].rs;
                l = mid + 1;
            }
        }
        return node[u].data;
    }
};
\end{codeline}
\begin{warnbox}
只复制根指针，修改路径上的节点新建。
\end{warnbox}


\chapter{平衡树}

\section{Splay树}
\entrymeta{$O(\log n)$ 均摊}
\subsection{P3391 【模板】文艺平衡树}\label{p3391-ux6a21ux677fux6587ux827aux5e73ux8861ux6811}


\begin{codeline}
class Tag {
public:
    int64 add = 0;
    bool rev = false;
};

class Splay {
public:
    class Node {
    public:
        int son[2]{}, fa = 0, size = 0;
        int64 data = 0, sum = 0;
        Tag tag;
    };
    vector<Node> node;
    int root = 0, n = 0;

    void push_up(int u) {
        int l = node[u].son[0], r = node[u].son[1];
        node[u].size = node[l].size + node[r].size + 1;
        node[u].sum = node[l].sum + node[r].sum + node[u].data;
    }
    void apply(int u, const Tag& t) {
        if(!u) {
            return;
        }
        node[u].data += t.add;
        node[u].sum += t.add * node[u].size;
        node[u].tag.add += t.add;
        if(t.rev) {
            swap(node[u].son[0], node[u].son[1]);
            node[u].tag.rev ^= true;
        }
    }
    void push_down(int u) {
        apply(node[u].son[0], node[u].tag);
        apply(node[u].son[1], node[u].tag);
        node[u].tag = Tag();
    }
    bool dir(int u) const {
        return node[node[u].fa].son[1] == u;
    }
    void rot(int u) {
        int fa = node[u].fa, g = node[fa].fa, d = dir(u);
        int v = node[u].son[d ^ 1];
        if(g) {
\end{codeline}
\vspace{-0.9\baselineskip}
\begin{codeline}[firstnumber=47]
            node[g].son[dir(fa)] = u;
        }
        node[u].fa = g;
        node[fa].son[d] = v;
        if(v) {
            node[v].fa = fa;
        }
        node[u].son[d ^ 1] = fa;
        node[fa].fa = u;
        push_up(fa);
        push_up(u);
    }
    void splay(int u, int goal = 0) {
        vector<int> st;
        for(int v = u; v != goal; v = node[v].fa) {
            st.push_back(v);
        }
        if(goal) {
            st.push_back(goal);
        }
        for(auto it = st.rbegin(); it != st.rend(); ++it) {
            push_down(*it);
        }
        while(node[u].fa != goal) {
            int fa = node[u].fa;
            if(node[fa].fa != goal) {
                rot(dir(u) == dir(fa) ? fa : u);
            }
            rot(u);
        }
        if(!goal) {
            root = u;
        }
    }
    /*
     * 由a[1..n]建立序列，两端补哨兵
     */
    int build(int l, int r, int fa, const vector<int64>& a) {
        if(l > r) {
            return 0;
        }
        int mid = (l + r) / 2, u = mid + 1;
        node[u].fa = fa;
        node[u].data = a[mid];
        node[u].son[0] = build(l, mid - 1, u, a);
        node[u].son[1] = build(mid + 1, r, u, a);
\end{codeline}
\vspace{-0.9\baselineskip}
\begin{codeline}[firstnumber=93]
        push_up(u);
        return u;
    }
    int kth(int k) {
        int u = root;
        while(true) {
            push_down(u);
            int s = node[node[u].son[0]].size;
            if(k <= s) {
                u = node[u].son[0];
            } else if(k == s + 1) {
                return u;
            } else {
                k -= s + 1;
                u = node[u].son[1];
            }
        }
    }
    int split(int l, int r) {
        int x = kth(l);
        splay(x);
        int y = kth(r + 2);
        splay(y, x);
        return node[y].son[0];
    }
    /*
     * 由a[1..n]建立序列，两端补哨兵
     */
    void build(const vector<int64>& a) {
        n = int(a.size()) - 1;
        vector<int64> b(n + 2);
        for(int i = 1; i <= n; i++) {
            b[i] = a[i];
        }
        node.assign(n + 3, Node());
        root = build(0, n + 1, 0, b);
    }
    /*
     * 对闭区间[l,r]施加加法和翻转标记t
     */
    void mod(int l, int r, const Tag& t) {
        apply(split(l, r), t);
        push_up(node[root].son[1]);
        push_up(root);
    }
    /*
\end{codeline}
\vspace{-0.9\baselineskip}
\begin{codeline}[firstnumber=139]
     * 翻转闭区间[l,r]
     */
    void rev(int l, int r) {
        mod(l, r, {0, true});
    }
    /*
     * 查询闭区间[l,r]的和
     */
    int64 query(int l, int r) {
        return node[split(l, r)].sum;
    }
    /*
     * 按当前序列顺序返回所有值，不含哨兵
     */
    vector<int64> values() {
        vector<int64> res;
        vector<int> st;
        int u = root, idx = 0;
        while(u || !st.empty()) {
            while(u) {
                push_down(u);
                st.push_back(u);
                u = node[u].son[0];
            }
            u = st.back();
            st.pop_back();
            if(++idx > 1 && idx <= n + 1) {
                res.push_back(node[u].data);
            }
            u = node[u].son[1];
        }
        return res;
    }
};
\end{codeline}
\begin{warnbox}
P3391 文艺平衡树：两端加哨兵，\texttt{kth} 返回的是带哨兵的下标。
\end{warnbox}

\section{PBDS 平衡树}
\entrymeta{$O(\log n)$}
\begin{codeline}
template <class Data = int>
class Splay {
public:
    class Node {
    public:
        int son[2]{}, fa = 0, cnt = 0, size = 0;
        Data data{};
    };
    vector<Node> node{Node()};
    int root = 0;

    bool dir(int u) const {
        return node[node[u].fa].son[1] == u;
    }
    void push_up(int u) {
        node[u].size = node[u].cnt + node[node[u].son[0]].size + node[node[u].son[1]].size;
    }
    void rot(int u) {
        int fa = node[u].fa, g = node[fa].fa, d = dir(u);
        int v = node[u].son[d ^ 1];
        if(g) {
            node[g].son[dir(fa)] = u;
        }
        node[u].fa = g;
        node[fa].son[d] = v;
        if(v) {
            node[v].fa = fa;
        }
        node[u].son[d ^ 1] = fa;
        node[fa].fa = u;
        push_up(fa);
        push_up(u);
    }
    void splay(int u) {
        while(node[u].fa) {
            int fa = node[u].fa;
            if(node[fa].fa) {
                rot(dir(u) == dir(fa) ? fa : u);
            }
            rot(u);
        }
        root = u;
    }
    int find(const Data& d) {
        int u = root, last = 0;
        while(u && !(node[u].data == d)) {
\end{codeline}
\vspace{-0.9\baselineskip}
\begin{codeline}[firstnumber=47]
            last = u;
            u = node[u].son[node[u].data < d];
        }
        if(u) {
            splay(u);
        } else if(last) {
            splay(last);
        }
        return u;
    }
    /*
     * 清空平衡树
     */
    void clear() {
        node.assign(1, Node());
        root = 0;
    }
    int size() const {
        return node[root].size;
    }
    /*
     * 插入d，重复值累计出现次数
     */
    void insert(const Data& d) {
        int u = root, fa = 0;
        while(u && !(node[u].data == d)) {
            fa = u;
            u = node[u].son[node[u].data < d];
        }
        if(!u) {
            u = int(node.size());
            node.emplace_back();
            node[u].data = d;
            node[u].fa = fa;
            if(fa) {
                node[fa].son[node[fa].data < d] = u;
            }
        }
        node[u].cnt++;
        push_up(u);
        splay(u);
    }
    /*
     * 删除一个d，返回是否存在并删除
     */
    bool remove(const Data& d) {
\end{codeline}
\vspace{-0.9\baselineskip}
\begin{codeline}[firstnumber=93]
        int u = find(d);
        if(!u) {
            return false;
        }
        if(--node[u].cnt) {
            push_up(u);
            return true;
        }
        int l = node[u].son[0], r = node[u].son[1];
        if(l) {
            node[l].fa = 0;
        }
        if(r) {
            node[r].fa = 0;
        }
        root = l;
        if(!l) {
            root = r;
        } else {
            int v = l;
            while(node[v].son[1]) {
                v = node[v].son[1];
            }
            splay(v);
            node[v].son[1] = r;
            if(r) {
                node[r].fa = v;
            }
            push_up(v);
        }
        return true;
    }
    /*
     * 返回d的排名，即小于d的元素个数加1
     */
    int rank(const Data& d) {
        find(d);
        if(!root) {
            return 1;
        }
        return node[node[root].son[0]].size + (node[root].data < d ? node[root].cnt : 0) + 1;
    }
    /*
     * 将第k小的值写入res，排名不存在时返回false
     */
    bool kth(int k, Data& res) {
\end{codeline}
\vspace{-0.9\baselineskip}
\begin{codeline}[firstnumber=139]
        if(k < 1 || k > size()) {
            return false;
        }
        int u = root;
        while(true) {
            int s = node[node[u].son[0]].size;
            if(k <= s) {
                u = node[u].son[0];
            } else if(k <= s + node[u].cnt) {
                break;
            } else {
                k -= s + node[u].cnt;
                u = node[u].son[1];
            }
        }
        res = node[u].data;
        splay(u);
        return true;
    }
    /*
     * 将严格小于d的最大值写入res，不存在时返回false
     */
    bool prev(const Data& d, Data& res) {
        return kth(rank(d) - 1, res);
    }
    /*
     * 将严格大于d的最小值写入res，不存在时返回false
     */
    bool next(const Data& d, Data& res) {
        int k = rank(d);
        if(root && node[root].data == d) {
            k += node[root].cnt;
        }
        return kth(k, res);
    }
};
\end{codeline}


\subsection{PBDS 实现}\label{pbds-ux5b9eux73b0}


\begin{codeline}
/*
 * PBDS 实现：用 __gnu_pbds 的红黑树，通过 TreeUpdate 策略维护子树大小，
 * 从而支持 rank / kth，省去手写旋转。
 */
using namespace __gnu_pbds;

class Node {
public:
    int k;
    mutable int cnt;

    bool operator<(const Node& o) const {
        return k < o.k;
    }
};

class MetaData {
public:
    int size;
};

MetaData METADATA_EP = {0};

template <class cIt, class It, class Cmp, class _A>
class TreeUpdate {
public:
    typedef MetaData metadata_type;

#define nd(x) ((x).m_p_nd)
#define lson(x) (nd(x)->m_p_left)
#define rson(x) (nd(x)->m_p_right)
#define md(x) ((x) == node_end() ? METADATA_EP : (metadata_type)(nd(x)->m_metadata))
#define rmd(x) ((metadata_type&)(nd(x)->m_metadata))

    virtual cIt node_begin() const = 0;

    virtual cIt node_end() const = 0;

    void operator()(It u, cIt end) {
        It ls = lson(u);
        It rs = rson(u);
        metadata_type data = {0};
        if(ls != end) {
            data.size += md(ls).size;
        }
        if(rs != end) {
\end{codeline}
\vspace{-0.9\baselineskip}
\begin{codeline}[firstnumber=47]
            data.size += md(rs).size;
        }
        data.size += (*u)->cnt;

        rmd(u) = data;
    }

    /*
     * 返回第 k 小的值
     */
    int kth(int k) {
        cIt u = node_begin();
        cIt end = node_end();

        while(u != end) {
            It ls = lson(u);
            It rs = rson(u);

            if(k <= md(ls).size) {
                u = ls;
            } else if(k <= md(ls).size + (*u)->cnt) {
                return (*u)->k;
            } else {
                k -= md(ls).size + (*u)->cnt;
                u = rs;
            }
        }
        return -1;
    }

    /*
     * 返回小于 k 的元素个数加一
     */
    int rank(int k) {
        cIt u = node_begin();
        cIt end = node_end();
        int rk = 1;

        while(u != end) {
            It ls = lson(u);
            It rs = rson(u);

            if(k <= (*u)->k) {
                u = ls;
            } else {
                rk += (*u)->cnt + md(ls).size;
\end{codeline}
\vspace{-0.9\baselineskip}
\begin{codeline}[firstnumber=93]
                u = rs;
            }
        }
        return rk;
    }

#undef nd
#undef lson
#undef rson
#undef md
#undef rmd
};

class BT : public tree<Node, null_type, less<Node>, rb_tree_tag, TreeUpdate> {
private:
    using Base = tree<Node, null_type, less<Node>, rb_tree_tag, TreeUpdate>;
    typedef Base::iterator It;

    void update(It node) {
        update_to_top(node.m_p_nd, this);
    }

public:
    /*
     * 插入 k，重复值累计出现次数
     */
    void insert(int k) {
        It it = lower_bound({k, 0});
        if(it != end() && it->k == k) {
            ++(it->cnt);
            update(it);
        } else {
            Base::insert({k, 1});
        }
    }

    /*
     * 删除一个 k
     */
    void erase(int k) {
        auto it = lower_bound(Node{k, 0});
        if(it != end() && it->cnt > 1) {
            --(it->cnt);
            update(it);
        } else {
            Base::erase(it);
\end{codeline}
\vspace{-0.9\baselineskip}
\begin{codeline}[firstnumber=139]
        }
    }
};

BT bt;

int main() {
    ios::sync_with_stdio(false);
    cin.tie(nullptr);
    int n, m, num;
    cin >> n >> m;

    for(int i = 1; i <= n; i++) {
        cin >> num;
        bt.insert(num);
    }

    int op, x, last = 0, ans = 0;
    while(m--) {
        cin >> op >> x;
        x ^= last;
        if(op == 1) {
            bt.insert(x);
        } else if(op == 2) {
            bt.erase(x);
        } else if(op == 3) {
            last = bt.rank(x);
            ans ^= last;
        } else if(op == 4) {
            last = bt.kth(x);
            ans ^= last;
        } else if(op == 5) {
            last = (--bt.lower_bound({x, 0}))->k;
            ans ^= last;
        } else if(op == 6) {
            last = (bt.upper_bound({x, 0}))->k;
            ans ^= last;
        }
    }
    cout << ans << '\n';
}
\end{codeline}
\begin{warnbox}
PBDS 版本用 \texttt{TreeNode} 维护子树 size；\texttt{erase} 前必须确认元素存在。
\end{warnbox}

\section{笛卡尔树}
\entrymeta{$O(n)$}
\begin{codeline}
const int N = 10000005;

class Node {
public:
    int x;
    int l, r;
};

Node node[N];

int main() {
    ios::sync_with_stdio(false);
    cin.tie(nullptr);
    cout.tie(nullptr);

    int n;
    cin >> n;
    stack<int> st;
    long long ansl = 0, ansr = 0;
    for(int i = 1; i <= n; i++) {
        cin >> node[i].x;
        int lst = 0;
        while(!st.empty() && node[st.top()].x > node[i].x) {
            lst = st.top();
            ansl ^= lst * (node[lst].l + 1ll);
            ansr ^= lst * (node[lst].r + 1ll);
            st.pop();
        }
        if(!st.empty()) {
            node[st.top()].r = i;
        }
        node[i].l = lst;
        st.push(i);
    }
    while(!st.empty()) {
        ansl ^= st.top() * (node[st.top()].l + 1ll);
        ansr ^= st.top() * (node[st.top()].r + 1ll);
        st.pop();
    }
    cout << ansl << ' ' << ansr << endl;
}
\end{codeline}
\begin{warnbox}
单调栈建树；下标从 1 开始。
\end{warnbox}


\chapter{分块}

\section{块状数组}
\entrymeta{$O(\sqrt n)$ 单次}
\begin{codeline}
const int N = 200005;
class BlockArray {
private:
    long long arr[N];
    long long idx[N];
    long long tag[500];
    long long n, S, bn;

    void bsort(long long b) {
        sort(idx + b * S,
             min(idx + (b + 1) * S, idx + n + 1),
             [this](long long A, long long B) {
                 return arr[A] < arr[B];
             });
    }

    long long bget(long long b, long long x) {
        return lower_bound(idx + b * S,
                           min(idx + (b + 1) * S, idx + n + 1),
                           x,
                           [this](long long A, long long B) {
                               return arr[A] < B;
                           }) -
               (idx + b * S);
    }

public:
    /*
     * 读入size个元素并初始化分块
     */
    void build(int size) {
        n = size;
        S = max(1ll, (long long)sqrt(n));
        bn = n / S + 1;
        for(long long i = 1; i <= n; i++) {
            cin >> arr[i];
            idx[i] = i;
        }

        for(long long i = 0; i < bn; i++) {
            bsort(i);
        }
    }

    /*
     * 将闭区间[l,r]的元素加x
\end{codeline}
\vspace{-0.9\baselineskip}
\begin{codeline}[firstnumber=47]
     */
    void add(long long l, long long r, long long x) {
        long long lb = l / S;
        long long rb = r / S;

        if(lb == rb) {
            for(long long i = l; i <= r; i++) {
                arr[i] += x;
            }
            bsort(lb);
        } else {
            for(long long i = l; i < (lb + 1) * S; i++) {
                arr[i] += x;
            }
            bsort(lb);
            for(long long i = lb + 1; i < rb; i++) {
                tag[i] += x;
            }
            for(long long i = rb * S; i <= r; i++) {
                arr[i] += x;
            }
            bsort(rb);
        }
    }

    /*
     * 统计闭区间[l,r]中小于x的元素个数
     */
    long long count(long long l, long long r, long long x) {
        long long lb = l / S;
        long long rb = r / S;
        long long cnt = 0;

        if(lb == rb) {
            for(long long i = l; i <= r; i++) {
                cnt += arr[i] + tag[lb] < x;
            }
        } else {
            for(long long i = l; i < (lb + 1) * S; i++) {
                cnt += arr[i] + tag[lb] < x;
            }
            for(long long i = lb + 1; i < rb; i++) {
                cnt += bget(i, x - tag[i]);
            }
            for(long long i = rb * S; i <= r; i++) {
                cnt += arr[i] + tag[rb] < x;
\end{codeline}
\vspace{-0.9\baselineskip}
\begin{codeline}[firstnumber=93]
            }
        }
        return cnt;
    }

};

BlockArray ba;
int main() {
    ios::sync_with_stdio(false);
    cin.tie(nullptr);
    long long n, q;
    cin >> n;
    q = n;
    ba.build(n);
    long long op;
    long long l, r, x;
    while(q--) {
        cin >> op;
        cin >> l;
        cin >> r;
        cin >> x;

        if(op == 0) {
            ba.add(l, r, x);
        }
        if(op == 1) {
            cout << (ba.count(l, r, x * x)) << '\n';
        }
    }
}
\end{codeline}
\begin{codeline}
const int N = 300005;
const int SQRTN = 600;
const long long INF = 1LL << 60;
class BlockArray {
private:
    int arr[N];
    long long eq[SQRTN];
    long long tag[SQRTN];
    int n, S, bn;

    void push_down(int b) {
        if(tag[b] != INF) {
            int r = min(n, (b + 1) * S - 1);
            for(int j = b * S; j <= r; j++) {
                arr[j] = tag[b];
            }
            tag[b] = INF;
        }
    }

    int bget(int l, int r, int x) {
        int lb = l / S, rb = r / S, ans = 0;
        if(lb == rb) {
            if(eq[lb] == x) {
                return r - l + 1;
            } else if(eq[lb] != INF) {
                if(r - l + 1 == S) {
                    eq[lb] = tag[lb] = x;
                } else {
                    push_down(lb);
                    eq[lb] = tag[lb] = INF;
                    for(int i = l; i <= r; i++) {
                        arr[i] = x;
                    }
                }
                return 0;
            } else {
                for(int i = l; i <= r; i++) {
                    ans += arr[i] == x;
                    arr[i] = x;
                }
                eq[lb] = r - l + 1 == S ? x : INF;
                tag[lb] = INF;
                return ans;
            }
        } else {
\end{codeline}
\vspace{-0.9\baselineskip}
\begin{codeline}[firstnumber=47]
            return 0;
        }
    }

public:
    /*
     * 读入size个元素并初始化分块
     */
    void build(int size) {
        n = size;
        S = max(1, (int)sqrt(n));
        bn = n / S + 1;
        for(int i = 1; i <= n; i++) {
            cin >> arr[i];
        }
        arr[0] = 0;
        for(int i = 0; i < bn; i++) {
            eq[i] = arr[i * S];
            tag[i] = INF;
        }
        eq[bn] = tag[bn] = INF;
        for(int i = 0; i <= n; i++) {
            if(arr[i] != eq[i / S]) {
                eq[i / S] = INF;
            }
        }
    }

    /*
     * 统计闭区间[l,r]中等于x的元素个数，并将该区间赋为x
     */
    int get(int l, int r, int x) {
        int lb = l / S, rb = r / S, ans = 0;
        if(lb == rb) {
            return bget(l, r, x);
        } else {
            ans += bget(l, (lb + 1) * S - 1, x);
            for(int i = lb + 1; i < rb; i++) {
                ans += bget(i * S, (i + 1) * S - 1, x);
            }
            ans += bget(rb * S, r, x);
            return ans;
        }
    }
};

\end{codeline}
\vspace{-0.9\baselineskip}
\begin{codeline}[firstnumber=93]
BlockArray ba;
int main() {
    ios::sync_with_stdio(false);
    cin.tie(nullptr);
    int n, q;
    cin >> n;
    q = n;
    ba.build(n);
    int op;
    int l, r, x;
    while(q--) {
        cin >> l;
        cin >> r;
        cin >> x;
        cout << (ba.get(l, r, x)) << '\n';
    }
}
\end{codeline}
\begin{codeline}
const int N = 100005;
const int SQRTN = 320;
const int INF = 0x7FFFFFFF;

int cnt[N];

class BlockArray {
private:
    int inv[N];
    int arr[N];
    int pre[SQRTN][N];
    int mode[SQRTN][SQRTN];
    int n, S, bn;

public:
    /*
     * 读入size个元素并初始化分块
     */
    void build(int size) {
        n = size;
        S = max(1, (int)sqrt(n));
        bn = n / S + 1;
        for(int i = 1; i <= n; i++) {
            cin >> arr[i];
            inv[i] = arr[i];
        }

        sort(inv + 1, inv + 1 + n);
        int len = unique(inv + 1, inv + 1 + n) - inv - 1;
        for(int i = 1; i <= n; i++) {
            arr[i] = lower_bound(inv + 1, inv + len + 1, arr[i]) - inv;
        }

        for(int b = 0; b < bn; b++) {
            if(b) {
                for(int i = 1; i <= len; i++) {
                    pre[b][i] = pre[b - 1][i];
                }
            }
            for(int i = b * S; i < min((b + 1) * S, n + 1); i++) {
                pre[b][arr[i]]++;
            }
        }

        for(int lb = 0; lb < bn; lb++) {
            int maxi = 0;
\end{codeline}
\vspace{-0.9\baselineskip}
\begin{codeline}[firstnumber=47]
            for(int rb = lb; rb < bn; rb++) {
                if(rb > lb) {
                    mode[lb][rb] = mode[lb][rb - 1];
                }
                for(int i = rb * S; i < min((rb + 1) * S, n + 1); i++) {
                    int v = arr[i];
                    int ans = pre[rb][v];
                    if(lb) {
                        ans -= pre[lb - 1][v];
                    }
                    if(ans > maxi || ans == maxi && v < mode[lb][rb]) {
                        mode[lb][rb] = v;
                        maxi = ans;
                    }
                }
            }
        }
    }

    /*
     * 查询闭区间[l,r]的众数，出现次数相同取较小值
     */
    int get(int l, int r) {
        int lb = l / S, rb = r / S;
        int maxi = 0;
        int ans = 0;
        if(lb == rb) {
            for(int i = l; i <= r; i++) {
                cnt[arr[i]] = 0;
            }
            for(int i = l; i <= r; i++) {
                ++cnt[arr[i]];
                if(cnt[arr[i]] > maxi || cnt[arr[i]] == maxi && arr[i] < ans) {
                    maxi = cnt[arr[i]];
                    ans = arr[i];
                }
            }
        } else {
            for(int i = l; i <= (lb + 1) * S; i++) {
                cnt[arr[i]] = pre[rb - 1][arr[i]] - pre[lb][arr[i]];
            }
            for(int i = rb * S; i <= r; i++) {
                cnt[arr[i]] = pre[rb - 1][arr[i]] - pre[lb][arr[i]];
            }
            maxi = cnt[mode[lb + 1][rb - 1]] = pre[rb - 1][mode[lb + 1][rb - 1]] - pre[lb][mode[lb + 1][rb - 1]];
            ans = mode[lb + 1][rb - 1];
\end{codeline}
\vspace{-0.9\baselineskip}
\begin{codeline}[firstnumber=93]
            for(int i = l; i < (lb + 1) * S; i++) {
                ++cnt[arr[i]];
                if(cnt[arr[i]] > maxi || cnt[arr[i]] == maxi && arr[i] < ans) {
                    maxi = cnt[arr[i]];
                    ans = arr[i];
                }
            }
            for(int i = rb * S; i <= r; i++) {
                ++cnt[arr[i]];
                if(cnt[arr[i]] > maxi || cnt[arr[i]] == maxi && arr[i] < ans) {
                    maxi = cnt[arr[i]];
                    ans = arr[i];
                }
            }
        }
        return inv[ans];
    }
};

BlockArray ba;
int main() {
    ios::sync_with_stdio(false);
    cin.tie(nullptr);
    int n, q;
    cin >> n;
    q = n;
    ba.build(n);
    int l, r;
    while(q--) {
        cin >> l;
        cin >> r;
        cout << (ba.get(l, r)) << '\n';
    }
}
\end{codeline}
\begin{warnbox}
块长 \texttt{S = max(1, sqrt(n))}；区间众数版本的 \texttt{pre} 按 n ≤ 1e5 开。
\end{warnbox}

\section{树分块}
\entrymeta{$O(\sqrt n)$}
\begin{codeline}
const int W = 64;
const int N = 40005;
const int LOGN = 17;
int cvt[N];

class Tree {
public:
    int S;
    class Node {
    public:
        int c, d;
        int p;
        int fa[LOGN];
        vector<int> oe;
    };

    class VNode {
    public:
        int fa;
        vector<int> oe;
        vector<bitset<N>> bm;
    };

    Node node[N];
    VNode vnode[N];

    /*
     * 添加无向边(u,v)
     */
    void add(int u, int v) {
        node[u].oe.push_back(v);
        node[v].oe.push_back(u);
    }

    /*
     * 以1为根建立树分块和倍增信息
     */
    void build(int n) {
        S = sqrt(n) / W * W;
        dfs1(1);
        for(int i = 1; i < LOGN; i++) {
            for(int j = 1; j <= n; j++) {
                node[j].fa[i] = node[node[j].fa[i - 1]].fa[i - 1];
            }
        }
        dfs2(1);
\end{codeline}
\vspace{-0.9\baselineskip}
\begin{codeline}[firstnumber=47]
    }

    /*
     * 返回x、y的最近公共祖先
     */
    int getLCA(int x, int y) {
        if(node[x].d < node[y].d) {
            swap(x, y);
        }
        for(int i = LOGN - 1; i >= 0 && node[x].d > node[y].d; i--) {
            if(node[x].d >= node[y].d + (1 << i)) {
                x = node[x].fa[i];
            }
        }
        if(x == y) {
            return x;
        } else {
            for(int i = LOGN - 1; i >= 0; i--) {
                if(node[x].fa[i] != node[y].fa[i]) {
                    x = node[x].fa[i];
                    y = node[y].fa[i];
                }
            }
            return node[x].fa[0];
        }
    }

    void dfs1(int r) {
        node[r].d = 1;
        stack<int> st;
        st.push(r);
        while(!st.empty()) {
            int u = st.top(), p = node[u].p;
            st.pop();
            if(node[u].d % S == 1) {
                vnode[p].oe.push_back(u);
                vnode[u].fa = p;
                p = node[u].p = u;
            }
            for(int v : node[u].oe) {
                if(v != node[u].fa[0]) {
                    node[v].fa[0] = u;
                    node[v].d = node[u].d + 1;
                    node[v].p = p;
                    st.push(v);
                }
\end{codeline}
\vspace{-0.9\baselineskip}
\begin{codeline}[firstnumber=93]
            }
        }
    }

    void dfs2(int r) {
        stack<int> st;
        st.push(r);
        while(!st.empty()) {
            int u = st.top(), fa = u;
            st.pop();
            vnode[u].bm.emplace_back();
            do {
                fa = node[fa].fa[0];
                vnode[u].bm[0][node[fa].c] = true;
            } while(fa != vnode[u].fa);
            for(auto& bm : vnode[vnode[u].fa].bm) {
                vnode[u].bm.emplace_back(vnode[u].bm[0] | bm);
            }
            for(int v : vnode[u].oe) {
                st.push(v);
            }
        }
    }

    /*
     * 查询x到y路径上的不同颜色数
     */
    int query(int x, int y) {
        int lca = getLCA(x, y);
        bitset<N> ans;
        for(int i = 1; i <= 2; i++, swap(x, y)) {
            if(node[lca].p != node[x].p) {
                int x0 = node[x].p;
                while(x != node[x].p && node[x].d > node[lca].d) {
                    ans[node[x].c] = true;
                    x = node[x].fa[0];
                }
                ans[node[x].c] = true;
                while((node[vnode[x].fa].d >= node[lca].d)) {
                    x = vnode[x].fa;
                }
                if(x0 != x) {
                    ans |= vnode[x0].bm[(node[x0].d - node[x].d) / S - 1];
                }
            }
            while(x != lca) {
\end{codeline}
\vspace{-0.9\baselineskip}
\begin{codeline}[firstnumber=139]
                ans[node[x].c] = true;
                x = node[x].fa[0];
            }
        }
        ans[node[lca].c] = true;
        return ans.count();
    }
};

Tree tree;

int main() {
    ios::sync_with_stdio(false);
    cin.tie(nullptr);
    int n, m, c;
    cin >> n;
    cin >> m;

    for(int i = 1; i <= n; i++) {
        cin >> c;
        cvt[i] = c;
        tree.node[i].c = c;
    }
    sort(cvt + 1, cvt + n + 1);
    int C = unique(cvt + 1, cvt + n + 1) - cvt - 1;
    for(int i = 1; i <= n; i++) {
        tree.node[i].c = lower_bound(cvt + 1, cvt + 1 + C, tree.node[i].c) - cvt;
    }

    int u, v;
    for(int i = 1; i < n; i++) {
        cin >> u;
        cin >> v;
        tree.add(u, v);
    }

    tree.build(n);
    int lastans = 0;
    for(int i = 1; i <= m; i++) {
        cin >> u;
        cin >> v;
        u ^= lastans;
        lastans = tree.query(u, v);
        cout << (lastans) << '\n';
    }
}
\end{codeline}
\begin{warnbox}
按深度分块，块内维护到祖先的 bitset。
\end{warnbox}


\chapter{动态树}

\section{Link-Cut Tree}
\entrymeta{$O(\log n)$ 均摊}
Link-Cut Tree 用一组 Splay 辅助树维护森林。下面的通用实现封装了核心操作， 题目只需给出 \texttt{Data}（聚合信息与 \texttt{operator+}）；需要额外维护虚子树等附加信息时， 再重写 \texttt{on\_access} / \texttt{on\_link} 两个钩子。


\begin{codeline}
const int N = 100005;

class Tag {
public:
    bool rev;
};

const Tag TAG_EP = {false};
const Tag TAG_REV = {true};

/*
 * 通用 LCT。Data 需提供 operator+ 和默认构造（单位元）；
 * self 保存节点自身的贡献，data = 左子 + self + 右子。
 * 需要维护虚子树等附加信息时，重写 on_access / on_link。
 */
template <class Data>
class LinkCutTree {
public:
    static const int MAXN = N;

    class Node {
    public:
        int son[2];
        int fa;
        Data data;
        Data self;
        Tag tag;

        void operator+=(const Tag& t) {
            if(t.rev) {
                swap(son[0], son[1]);
                tag.rev ^= t.rev;
            }
        }
    };

    Node node[MAXN];

    /*
     * 实链切换时的附加维护，默认无操作
     */
    virtual void on_access(int, int, int) {
    }

    /*
     * 加边后的附加维护，默认无操作
\end{codeline}
\vspace{-0.9\baselineskip}
\begin{codeline}[firstnumber=47]
     */
    virtual void on_link(int, int) {
    }

#define ls(u) (node[u].son[0])
#define rs(u) (node[u].son[1])
#define fa(u) (node[u].fa)
#define data(u) (node[u].data)
#define tag(u) (node[u].tag)

    bool dir(const int u) {
        return rs(fa(u)) == u;
    }

    bool is_root(const int u) {
        return ls(fa(u)) != u && rs(fa(u)) != u;
    }

    void push_up(const int u) {
        data(u) = data(ls(u)) + node[u].self + data(rs(u));
    }

    void push_down(const int u) {
        if(ls(u)) {
            node[ls(u)] += tag(u);
        }
        if(rs(u)) {
            node[rs(u)] += tag(u);
        }
        tag(u) = TAG_EP;
    }

    void update(int u) {
        vector<int> st{u};
        for(int v = u; !is_root(v);) {
            v = fa(v);
            st.push_back(v);
        }
        for(auto it = st.rbegin(); it != st.rend(); ++it) {
            push_down(*it);
        }
    }

    void rot(int u) {
        int y = fa(u), z = fa(y), k = dir(u), v = node[u].son[!k];
        if(!is_root(y)) {
\end{codeline}
\vspace{-0.9\baselineskip}
\begin{codeline}[firstnumber=93]
            node[z].son[dir(y)] = u;
        }
        node[y].son[k] = v;
        node[u].son[!k] = y;
        if(v) {
            fa(v) = y;
        }
        fa(y) = u;
        fa(u) = z;
        push_up(y);
        push_up(u);
    }

    void splay(int u) {
        update(u);
        for(int f = fa(u); !is_root(u); rot(u), f = fa(u)) {
            if(!is_root(f)) {
                rot(dir(u) == dir(f) ? f : u);
            }
        }
    }

    /*
     * 打通原树根到 u 的实链
     */
    int access(int u) {
        int v = 0;
        for(; u; v = u, u = fa(u)) {
            splay(u);
            on_access(u, rs(u), v);
            rs(u) = v;
            push_up(u);
        }
        return v;
    }

    /*
     * 将 u 设为所在原树的根
     */
    void make_root(int u) {
        access(u);
        splay(u);
        node[u] += TAG_REV;
    }

    /*
\end{codeline}
\vspace{-0.9\baselineskip}
\begin{codeline}[firstnumber=139]
     * 返回 u 所在原树的根
     */
    int find_root(int u) {
        access(u);
        splay(u);
        push_down(u);
        while(ls(u)) {
            u = ls(u);
            push_down(u);
        }
        splay(u);
        return u;
    }

    /*
     * 将 x 到 y 的路径暴露到以 y 为根的辅助树，x、y 须连通
     */
    void split(int x, int y) {
        make_root(x);
        access(y);
        splay(y);
    }

    /*
     * 连接 x、y；只在两点不连通时加边
     */
    void link(int x, int y) {
        make_root(x);
        if(find_root(y) != x) {
            splay(y);
            fa(x) = y;
            on_link(x, y);
            push_up(y);
        }
    }

    /*
     * 断开 x、y 之间的直接边
     */
    void cut(int x, int y) {
        make_root(x);
        if(find_root(y) == x && fa(y) == x && !ls(y)) {
            fa(y) = rs(x) = 0;
            push_up(x);
        }
    }
\end{codeline}
\vspace{-0.9\baselineskip}
\begin{codeline}[firstnumber=185]

#undef ls
#undef rs
#undef fa
#undef data
#undef tag
};
\end{codeline}


\subsection{P3690 【模板】动态树（LCT）}\label{p3690-ux6a21ux677fux52a8ux6001ux6811lct}


\begin{codeline}
class Data {
public:
    int sum;

    Data operator+(const Data& other) const {
        return {sum ^ other.sum};
    }
};

int n, m;

LinkCutTree<Data> lct;

int main() {
    ios::sync_with_stdio(false);
    cin.tie(nullptr);
    cin >> n;
    cin >> m;

    for(int i = 1; i <= n; i++) {
        cin >> lct.node[i].self.sum;
        lct.node[i].data.sum = lct.node[i].self.sum;
    }

    int op, x, y;
    while(m--) {
        cin >> op;
        cin >> x;
        cin >> y;
        if(op == 0) {
            lct.split(x, y);
            cout << (lct.node[y].data.sum) << '\n';
        } else if(op == 1) {
            lct.link(x, y);
        } else if(op == 2) {
            lct.cut(x, y);
        } else if(op == 3) {
            lct.make_root(x);
            lct.node[x].self.sum = y;
            lct.push_up(x);
        }
    }
}
\end{codeline}
\begin{warnbox}
\texttt{Data::operator+} 需可交换（反转只交换儿子，不重算自身）；动态重心需重写 \texttt{on\_access} / \texttt{on\_link}。
\end{warnbox}


\chapter{其他结构}

\section{KD-Tree}
\entrymeta{$O(\sqrt n)$ 期望}
\begin{codeline}
const int64 INF = 0x3f3f3f3f3f3f3f3f;
const int N = 400005;
const int KD = 2;
int64 n;

class Point {
public:
    int64 coord[KD];
};
Point point[N];

class Node {
public:
    int l, r;
    int64 coord[KD];
    int64 min_coord[KD];
    int64 max_coord[KD];
};

bool cmp0(const Point& A, const Point& B) {
    return A.coord[0] < B.coord[0];
}

bool cmp1(const Point& A, const Point& B) {
    return A.coord[1] < B.coord[1];
}

class KDTree {
public:
    static const int MAXN = 400005;
    int kd = KD;
    int root = -1;
    int node_cnt = 0; // 0 for root;
    Node node[MAXN];

    int64 get_distance_squared(int tgt, int cur) {
        if(cur == -1) {
            return INF;
        }
        int64 gap, d2 = 0.;
        for(int i = 0; i < kd; i++) {
            gap = node[tgt].coord[i] - node[cur].coord[i];
            d2 += gap * gap;
        }
        return d2;
    }
\end{codeline}
\vspace{-0.9\baselineskip}
\begin{codeline}[firstnumber=47]

    int64 potential_min_dis_squared(int tgt, int cur) {
        if(cur == -1) {
            return INF;
        }
        int64 gap, d2 = 0.;
        for(int i = 0; i < kd; i++) {
            if(node[tgt].coord[i] < node[cur].min_coord[i]) {
                gap = node[cur].min_coord[i] - node[tgt].coord[i];
            } else if(node[tgt].coord[i] > node[cur].max_coord[i]) {
                gap = node[tgt].coord[i] - node[cur].max_coord[i];
            } else {
                gap = 0;
            }
            d2 += gap * gap;
        }
        return d2;
    }

    void init_node(const Point& point, int cur) {
        for(int i = 0; i < kd; i++) {
            node[cur].coord[i] =
                node[cur].min_coord[i] =
                    node[cur].max_coord[i] = point.coord[i];
        }
        node[cur].l = -1;
        node[cur].r = -1;
    }

    void push_up(int cur) {
        for(int i = 0; i < kd; i++) {
            node[cur].min_coord[i] = node[cur].coord[i];
            node[cur].max_coord[i] = node[cur].coord[i];

            if(node[cur].l != -1) {
                node[cur].min_coord[i] = min(node[cur].min_coord[i], node[node[cur].l].min_coord[i]);
                node[cur].max_coord[i] = max(node[cur].max_coord[i], node[node[cur].l].max_coord[i]);
            }

            if(node[cur].r != -1) {
                node[cur].min_coord[i] = min(node[cur].min_coord[i], node[node[cur].r].min_coord[i]);
                node[cur].max_coord[i] = max(node[cur].max_coord[i], node[node[cur].r].max_coord[i]);
            }
        }
    }

\end{codeline}
\vspace{-0.9\baselineskip}
\begin{codeline}[firstnumber=93]
    /*
     * 查询所有点对的最小距离平方
     */
    int64 batch_query() {
        int64 cur_ans = INF;
        int cur;
        for(int tgt = 0; tgt < node_cnt; tgt++) {
            stack<int> s;
            s.push(root);

            while(!s.empty()) {
                cur = s.top();
                s.pop();
                if(tgt == cur || potential_min_dis_squared(tgt, cur) > cur_ans) {
                    continue;
                } else {
                    cur_ans = min(cur_ans, get_distance_squared(tgt, cur));
                    if(node[cur].l != -1) {
                        s.push(node[cur].l);
                    }
                    if(node[cur].r != -1) {
                        s.push(node[cur].r);
                    }
                }
            }
        }
        return cur_ans;
    }

    /*
     * 按维度轮流划分point[range_l..range_r]，返回根节点编号
     */
    int build(int range_l, int range_r, int d) {
        if(range_l == range_r) {
            init_node(point[range_l], node_cnt);
            return node_cnt++;
        } else if(range_l > range_r) {
            return -1;
        } else {
            int node_idx = node_cnt++;
            int mid = (range_r + range_l) >> 1;
            int d_idx = d % 2;

            if(d_idx == 0) {
                nth_element(point + range_l, point + mid, point + range_r, cmp0);
            } else {
\end{codeline}
\vspace{-0.9\baselineskip}
\begin{codeline}[firstnumber=139]
                nth_element(point + range_l, point + mid, point + range_r, cmp1);
            }
            init_node(point[mid], node_idx);
            node[node_idx].l = build(range_l, mid - 1, d + 1);
            node[node_idx].r = build(mid + 1, range_r, d + 1);
            push_up(node_idx);
            return node_idx;
        }
    }
};

KDTree kdTree;

int main() {
    ios::sync_with_stdio(false);
    cin.tie(nullptr);
    cin >> n;
    for(int i = 0; i < n; i++) {
        for(int j = 0; j < KD; j++) {
            cin >> point[i].coord[j];
        }
    }

    kdTree.root = kdTree.build(0, int(n - 1), 0);

    cout << (kdTree.batch_query());
}
\end{codeline}
\begin{warnbox}
重构阈值与维度相关；最近点对需要剪枝。
\end{warnbox}

\section{可持久化字典树}
\entrymeta{$O(\log V)$}
MAXH为1..64，num和x为MAXH位无符号整数。查询排名不超过两个版本的计数差。


\begin{codeline}
template <int MAXH = 32>
class PersistentTrie {
public:
    class Node {
    public:
        int son[2]{};
        int cnt = 0;
    };
    int ncnt;
    vector<Node> node;
    vector<int> root;

    /*
     * 初始化，版本0为空Trie
     * n为预计插入次数，用于预留空间
     */
    void init(int n = 0) {
        ncnt = 0;
        node.assign(1, Node());
        node.reserve(size_t(n) * (MAXH + 1) + 1);
        root.assign(1, 0);
    }
    int getNode(int lst) {
        node.push_back(node[lst]);
        node[++ncnt].cnt++;
        return ncnt;
    }
    /*
     * 在版本ver上插入num，返回新版本号
     * 原版本保持不变
     */
    int push(int ver, uint64 num) {
        int u = getNode(root[ver]);
        root.push_back(u);
        for(int d = MAXH - 1; d >= 0; d--) {
            int c = (num >> d) & 1;
            int v = getNode(node[u].son[c]);
            node[u].son[c] = v;
            u = v;
        }
        return int(root.size()) - 1;
    }
    /*
     * 版本vr减去版本vl后的Trie，与x异或后第k大的原数
     * 左开右闭，vl为vr的祖先版本，k从1开始
     */
\end{codeline}
\vspace{-0.9\baselineskip}
\begin{codeline}[firstnumber=47]
    uint64 getKth(int vl, int vr, uint64 x, int k = 1) const {
        int l = root[vl], r = root[vr];
        uint64 res = 0;
        for(int d = MAXH - 1; d >= 0; d--) {
            int c = ((x >> d) & 1) ^ 1;
            int cnt = node[node[r].son[c]].cnt - node[node[l].son[c]].cnt;
            if(k > cnt) {
                k -= cnt;
                c ^= 1;
            }
            res |= uint64(c) << d;
            l = node[l].son[c];
            r = node[r].son[c];
        }
        return res;
    }
};
\end{codeline}
\begin{warnbox}
版本区间查询用 \texttt{root[r] - root[l-1]}。
\end{warnbox}

\section{ST 表}
\entrymeta{$O(n\log n)$ 建表，$O(1)$ 查询}
\begin{codeline}
template <int n>
class SparseTable {
public:
    static constexpr int logn = 32 - __builtin_clz(n);
    int64 st[logn][n + 1]{};
    /*
     * 由st[0][1..n]预处理区间最大值
     */
    void calc() {
        for(int j = 1; j < logn; j++) {
            for(int i = 1; i + (1 << j) - 1 <= n; i++) {
                st[j][i] = max(st[j - 1][i], st[j - 1][i + (1 << (j - 1))]);
            }
        }
    }
    /*
     * 查询闭区间[l,r]的最大值
     */
    int64 query(int l, int r) const {
        int k = 31 - __builtin_clz(r - l + 1);
        return max(st[k][l], st[k][r - (1 << k) + 1]);
    }
};
\end{codeline}
\begin{warnbox}
\texttt{logn = 32 - \_\_builtin\_clz(n)}；要求 n ≥ 1。
\end{warnbox}

